\documentclass[11pt]{article}
\usepackage[letterpaper,margin=1in]{geometry}
\usepackage{amsmath,amssymb,amsthm,mathtools}
\usepackage{enumitem}
\usepackage{booktabs}
\usepackage{xurl}
\usepackage[hidelinks]{hyperref}
\usepackage{microtype}
\usepackage{seqsplit}

\newtheorem{theorem}{Theorem}[section]
\newtheorem{lemma}[theorem]{Lemma}
\newtheorem{proposition}[theorem]{Proposition}
\newtheorem{corollary}[theorem]{Corollary}
\theoremstyle{definition}
\newtheorem{definition}[theorem]{Definition}
\theoremstyle{remark}
\newtheorem{remark}[theorem]{Remark}

\newcommand{\E}{\mathbb{E}}
\newcommand{\Prob}{\mathbb{P}}
\newcommand{\R}{\mathbb{R}}
\newcommand{\N}{\mathbb{N}}
\newcommand{\fh}{\widehat f}
\newcommand{\cert}{\mathrm{CERT}}
\newcommand{\feas}{\mathcal Q}
\newcommand{\arrow}{\longrightarrow}
\makeatletter
\g@addto@macro{\UrlBreaks}{\do\_\do\.}
\makeatother
\urlstyle{tt}
\newcommand{\lean}[1]{\nolinkurl{#1}}
\newcommand{\repository}{\texttt{[repository URL to be inserted]}}
\hypersetup{
  pdftitle={A (3.49...) to the power k bound for diagonal Ramsey numbers},
  pdfauthor={Zhipu Zhao, William Booth-Clibborn, Casper Schwahn}
}

\title{A $(3.49\ldots)^k$ bound for diagonal Ramsey numbers}
\author{Zhipu Zhao \and William Booth-Clibborn \and Casper Schwahn}
\date{8 October 2026}

\begin{document}
\maketitle

\begin{abstract}
The diagonal Ramsey number $R(k,k)$ is the least order of a complete graph such that every red/blue colouring
of its edges contains a monochromatic clique of order $k$. We prove $R(k,k)\le b_*^k$ for all sufficiently large $k$,
where $b_*=3.49862379380751\ldots$ is an explicit constant. In particular, $R(k,k)\le3.4986238^k<3.5^k$ for all sufficiently
large $k$, improving the bound $3.69507^k$ of Lu and Wang. The classical bound $R(k,k)<4^k$ was proved by
Erd\H{o}s and Szekeres in 1935. Its exponential base was first reduced by Campos, Griffiths, Morris and
Sahasrabudhe in 2023. Our proof strengthens their algorithm by tracking a potential that gives greater
weight to vertices with high red density. An explicit integral certifies each use of the resulting local step.
We then use the improved off-diagonal bounds as inputs to further applications of the method.
Induction on the number of vertices justifies these dependencies without assuming the bounds to be proved.
The finite numerical inequalities are certified with exact arithmetic, and the complete argument, including
the off-diagonal input, has been formalised in Lean.
\end{abstract}

\setcounter{tocdepth}{1}
\tableofcontents
\clearpage

%======================================================================
\section{Introduction}
%======================================================================

For positive integers $k$ and $l$, the Ramsey number $R(k,l)$ is the least $N$ such that every colouring of the edges
of the complete graph $K_N$ in red and blue contains a red $K_k$ or a blue $K_l$. Erd\H{o}s and Szekeres~\cite{ES}
proved $R(k,k)\le\binom{2k-2}{k-1}<4^k$ in 1935. For almost ninety years the base $4$ was not improved. Campos,
Griffiths, Morris and Sahasrabudhe~\cite{CGMS} then showed $R(k,k)\le(4-\varepsilon)^k$ for a small fixed
$\varepsilon>0$, by a new algorithm that builds large monochromatic ``books''. Gupta, Ndiaye, Norin and Wei~\cite{GNNW}
optimised their method to reach about $3.8^k$. Lu and Wang~\cite{LW} recently obtained
\[
R(k,k)\le 3.69507^k\qquad\text{for all large }k,
\]
with a computer-assisted proof that is also checked in Lean.

Our main result is the following.

\begin{theorem}\label{thm:main}
Put
\[
z^*=\frac{626184844701}{500000000000}=1.252369689402,\qquad b_*=e^{z^*}.
\]
There is an integer $K\ge1$ such that, for every integer $k\ge K$,
\[
R(k,k)\le b_*^k\le
\left(\frac{34986238}{10000000}\right)^k<(7/2)^k.
\]
Here $b_*=3.49862379380751\ldots$ and $34986238/10000000=3.4986238$.
There is also an absolute constant $C\ge1$ such that $R(k,k)\le C\,e^{z^*k}$ for every $k\ge1$.
\end{theorem}

The notation $3.49\ldots$ in the title denotes $b_*$, not the smaller decimal $3.49$.
The proof gives a profile with exponent strictly below $z^*$. This strict gap lets us absorb the constant factor
and obtain the eventual bound with base $b_*$ itself. The decimal upper bound $3.4986238$ is certified separately.
The constant $C$ and the threshold $K$ are not effective, because
Theorem~\ref{thm:envelope} uses a compactness argument.

\paragraph{Formal verification and data.} We formalise the proof in Lean~4~\cite{Lean4} using
Mathlib~\cite{Mathlib}, Lu and Wang's development~\cite{LWcode}, and its dependency
\texttt{RamseyLean}~\cite{RamseyLean}. Section~\ref{sec:verification} explains the numerical checks, the modular
Lean build, and the final statement and axiom checks. The paper, certificate and reproduction instructions will
be available at \repository.

\paragraph{What is new.} The proof follows the book algorithm of Campos, Griffiths, Morris and
Sahasrabudhe~\cite{CGMS}, in the finite form developed by Lu and Wang~\cite{LW}. A monochromatic \emph{book}
consists of a clique, called its spine, and a set of further vertices, called its pages, each joined to every
spine vertex in the same colour. We change two parts of the analysis.

\emph{The local step.} The book algorithm keeps a pair of vertex sets $(X,Y)$ with many red edges between them. At
each step it either moves into the red neighbourhood of a vertex of $X$, or into its blue neighbourhood, or builds a
blue book. Previous analyses track the red edge density between $X$ and $Y$. We track instead the potential
\[
\Phi_c(X,Y)=|X|^w\,|Y|\,\Big(\frac1{|Y|}\sum_{y\in Y}(g_y-c)_+^{\,s}\Big)^{r/s},
\]
where $g_y$ is the proportion of $X$ joined in red to $y$, the shift $c$ is slightly below the target density, and
$w>0$, $r$ large, $s=\beta r$ with $0<\beta<1$, and $u_+=\max\{u,0\}$. Passing to subsets preserves the possible
clique outcomes, so we may choose a pair that maximises this potential. Maximality controls both the column
densities and the contribution of each row. When few rows have large blue degree,
Theorem~\ref{thm:envelope} rules out simultaneous failure of the red and blue tests if the explicit integral
$\cert(p,\mu,w,\beta,x)$ is negative. If many rows have large blue degree, we extract a blue book instead.
These alternatives give the local step used in the induction.

\emph{Feedback into the off-diagonal input.} The local step needs, as a stopping rule, an upper bound for $R(a,b)$
for all pairs $(a,b)$, not only the diagonal. Lu and Wang supply such a bound, $\log R(a,b)\le \fh(a,b)+o(a+b)$, through
an explicit function $\fh$. The bounds our algorithm produces are better than $\fh$ in a range of ratios $b/a$. We
therefore use these better bounds as off-diagonal inputs. Section~\ref{sec:banks} proves that a finite list of
such bounds may depend on each other, even cyclically, provided every dependency is applied to a strictly smaller
colouring. The proof is an induction on the number of vertices, not an assumption that the improved bounds already
hold. The finite collection used for Theorem~\ref{thm:main} contains all the bounds and comparisons needed to
justify this induction.

\paragraph{Outline.} Section~\ref{sec:prelim} fixes notation and states the input from Lu and Wang.
Section~\ref{sec:local} proves the local step. Section~\ref{sec:closure} turns it into a bound for colourings that
satisfy a stopping rule. Section~\ref{sec:banks} sets up profiles and banks and proves that a bank satisfying finitely
many inequalities is valid. Section~\ref{sec:certificate} describes the certified bank and completes the proof of
Theorem~\ref{thm:main}. Section~\ref{sec:verification} describes the verification.

%======================================================================
\section{Notation and the off-diagonal input}\label{sec:prelim}
%======================================================================

\subsection{Colourings}

A \emph{host} is a red/blue colouring of the edges of a complete graph on a finite vertex set $V$; we write $|V|$ for
its order. For disjoint finite sets $X,Y\subseteq V$, the red density between them is
$d_R(X,Y)=e_R(X,Y)/(|X||Y|)$, where $e_R(X,Y)$ is the number of red edges with one end in $X$ and one in $Y$. The red
density of $V$ itself is $d_R(V)=e_R(V)/\binom{|V|}2$, and $d_B(V)=1-d_R(V)$. For positive integers $a,b$, we write
$V\supseteq(a,b)$ if $V$ contains a red $K_a$ or a blue $K_b$, and $N\arrow(a,b)$ if every host on $N$ vertices does.
Thus $R(a,b)=\min\{N:N\arrow(a,b)\}$. The property $N\arrow(a,b)$ is monotone in $N$. We use the Erd\H{o}s--Szekeres
bound
\[
\binom{a+b-2}{a-1}\arrow(a,b).
\]
Exchanging the two colours turns a host with red density $d$ into one with red density $1-d$, and a red $K_a$ or
blue $K_b$ into a blue $K_a$ or red $K_b$. We call this \emph{renaming the colours}.

\begin{lemma}[balanced cut]\label{lem:cut}
Let $W$ be a host with $|W|\ge3$ and red density $d$. Then $W=X\sqcup Y$ with $d_R(X,Y)\ge d$ and
$|X|,|Y|\ge|W|/3$.
\end{lemma}

\begin{proof}
Average $d_R(X,W\setminus X)$ over all $X\subseteq W$ with $|X|=\lfloor|W|/2\rfloor$. Every pair of distinct vertices
is separated by the same number of these sets, so the average of $e_R(X,W\setminus X)$ is a fixed multiple of
$e_R(W)$, and the average of $|X||W\setminus X|$ is the same multiple of $\binom{|W|}2$. Hence some $X$ has
$d_R(X,W\setminus X)\ge d$. Both parts have at least $\lfloor|W|/2\rfloor\ge|W|/3$ vertices when $|W|\ge3$.
\end{proof}

\subsection{The Lu--Wang source function}

Lu and Wang~\cite{LW} define a function $f:[0,1]\to\R$ by
\[
f(t)=(1+t)\log(1+t)-t\log t+t\,e^{-t}P(t)\qquad(0\log0=0).
\]
Here $P$ is the piecewise cubic Hermite interpolant of $2280$ rational knots $0=t_0<\dots<t_{2279}=1$, with prescribed
rational values and derivatives at the knots. On each interval between consecutive knots, $P$ is the cubic
polynomial matching the values and first derivatives at both endpoints. Thus $P$ is $C^1$.
These data are the file \texttt{data/source.json} of the Lu--Wang
repository~\cite{LWcode} (SHA-256 \texttt{f3e4a4ea\ldots c08c}; Appendix~\ref{app:data}). The
\emph{symmetric extension} of $f$ is
\[
\fh(\lambda)=\begin{cases}f(\lambda),&0<\lambda\le1,\\ \lambda f(1/\lambda),&\lambda\ge1,\end{cases}\qquad
\fh(a,b)=a\,\fh(b/a)\quad(a,b>0).
\]
So $\fh(a,b)=\fh(b,a)$, and $\fh(a,b)$ is homogeneous of degree one.

\begin{theorem}[Lu--Wang uniform source bound {\cite[eq.~(32), Lemma B.1]{LW}}]\label{thm:LW}
For every $\varepsilon>0$ there is $C_\varepsilon\ge1$ such that, for all integers $a,b\ge1$ and $N$,
\[
N\ge C_\varepsilon\exp\big(\fh(a,b)+\varepsilon(a+b)\big)\quad\Longrightarrow\quad N\arrow(a,b).
\]
\end{theorem}

Lu and Wang state and prove this in Lean as \texttt{RamseyCurrent.FiniteSourceBound}. Their theorem quantifies over
all pairs, however unbalanced. Our Lean development imports their proof, not only their statement
(Section~\ref{sec:verification}).

\begin{lemma}[symmetric concavity]\label{lem:concave}
$\fh$ is concave on $(0,\infty)$. Consequently, for every $t>0$ with $t\ne1$, $\fh$ is differentiable at $t$ and
$\fh(s)\le\fh(t)+\fh'(t)(s-t)$ for all $s>0$. At $t=1$ the same inequality holds with any slope in
$[f(1)-f'(1),\,f'(1)]$.
\end{lemma}

\begin{proof}
Lu and Wang prove that $f$ is concave on $[0,1]$. For $\lambda>1$ the function $\lambda f(1/\lambda)$ is the
perspective of a concave function, so it is concave on $[1,\infty)$. A function that is concave on $(0,1]$ and on
$[1,\infty)$, and continuous at $1$, is concave on $(0,\infty)$ if its left derivative at $1$ is at least its right
derivative there. The left derivative is $f'(1)$; the right derivative is $f(1)-f'(1)$. So the condition is
$2f'(1)-f(1)\ge0$. Writing $v,d$ for the value and derivative of $P$ at $1$, one finds
$e\,(2f'(1)-f(1))=2d-v$, which is a positive rational number for the Lu--Wang data. The tangent bounds are the usual
supergradient inequalities of a concave function.
\end{proof}

The tangent bounds are what the certificate uses to compare $\fh$ with straight lines.

%======================================================================
\section{The local step}\label{sec:local}
%======================================================================

\subsection{A potential that weights high-density columns}

We need a quantity that measures both the sizes of the two sets and the red density seen by each vertex of $Y$.
The moment defined below gives more weight to vertices with unusually many red neighbours in $X$.
Maximising the resulting potential will give the row and column inequalities used in the local step.

Fix a host on $V$, and disjoint nonempty sets $X,Y\subseteq V$. For $v\in X$ and $y\in Y$ put $h_{vy}=1$ if $vy$ is red
and $0$ otherwise. The \emph{column density} of $y$ and the \emph{row density} of $v$ are
\[
g_y=\frac1{|X|}\sum_{v\in X}h_{vy},\qquad \pi_v=\frac1{|Y|}\sum_{y\in Y}h_{vy}.
\]
Both average to $d_R(X,Y)$. Fix reals $w>0$, $0<\beta<1$ and $r>w$ with $s=\beta r>1$, so $1<s<r$. For a real shift
$c$ let
\[
H_c(X,Y)=\Big(\frac1{|Y|}\sum_{y\in Y}(g_y-c)_+^{\,s}\Big)^{1/s},\qquad \Phi_c(X,Y)=|X|^w\,|Y|\,H_c(X,Y)^r,
\]
with $\Phi_c=0$ if $X$ or $Y$ is empty. We call $(X,Y)$ \emph{maximal} (for $c,w,r,s$) if $\Phi_c(X',Y)\le\Phi_c(X,Y)$
for every nonempty $X'\subseteq X$ and $\Phi_c(X,Y')\le\Phi_c(X,Y)$ for every nonempty $Y'\subseteq Y$.

The potential rewards large sets and large column densities. The moment of order $s$ weights the columns with high
red density more than an average does. The weight $w$ sets the relative value of a vertex of $X$ against a vertex of
$Y$.

Throughout this section we write $N=|X|$, $M=|Y|$, $u_y=g_y-c$, $H=H_c(X,Y)$, and $\E_X,\E_Y$ for uniform averages.

\begin{lemma}[columns at a maximal pair]\label{lem:columns}
Suppose $(X,Y)$ is maximal and $H>0$. Then $u_y\ge(1-\beta)^{1/s}H>0$ for every $y\in Y$. In particular
$\alpha:=\E_Yu_y\ge(1-\beta)^{1/s}H$.
\end{lemma}

\begin{proof}
For nonempty $Y'\subseteq Y$ we have $\Phi_c(X,Y')=N^w|Y'|^{1-r/s}\big(\sum_{y\in Y'}(u_y)_+^s\big)^{r/s}$, and
$1-r/s<0$. If some $u_{y_0}\le0$ and $M\ge2$, deleting $y_0$ keeps the sum and lowers $|Y'|$, which strictly increases
$\Phi_c$; this contradicts maximality. If $M=1$ then $u_y=H>0$. Now let $M\ge2$, fix $y$ and put $a=u_y/H$, so that
$u_y^s/\sum_{y'}u_{y'}^s=a^s/M$. Maximality at $Y\setminus\{y\}$, raised to the power $s/r$, gives
\[
1-\frac{a^s}M\le\Big(1-\frac1M\Big)^{1-s/r}\le1-\frac{1-\beta}M,
\]
by Bernoulli's inequality, since $1-s/r=1-\beta\in(0,1)$. Hence $a^s\ge1-\beta$.
\end{proof}

\begin{lemma}[rows at a maximal pair]\label{lem:rows}
Suppose $(X,Y)$ is maximal and $H>0$. For $v\in X$ put
\[
V_v=\frac{\E_Y\big[u_y^{s-1}(h_{vy}-c)\big]}{H^s}.
\]
Then $\E_XV_v=1$, and for every $z\in[0,1]^X$ with $t=\E_Xz_v>0$,
\begin{equation}\label{eq:support}
\E_X[z_vV_v]\le t^{1-w/r}.
\end{equation}
In particular $V_v\ge1-w/r$ for every $v\in X$.
\end{lemma}

The variable $z_v$ allows a fractional choice of rows. Inequality~\eqref{eq:support} bounds the total
contribution of such a choice by its size. The final assertion says that no row contributes much less than
the average contribution, which is $1$.

\begin{proof}
Averaging $h_{vy}-c$ over $v$ gives $u_y$, so $\E_XV_v=\E_Y[u_y^s]/H^s=1$. For $z\in[0,1]^X$ let
$b_y=\E_X[z_v(h_{vy}-c)]$ and define
\[
\Lambda(z)=t\,\Big(\frac{\|(b/t)_+\|_s}{H}\Big)^{r/w},\qquad \|v\|_s=\big(\E_Y|v_y|^s\big)^{1/s},
\]
with $\Lambda(0)=0$. The map $v\mapsto(\|v_+\|_s/H)^{r/w}$ is convex, because $v\mapsto\|v_+\|_s$ is convex and
nonnegative and $r/w>1$. So $\Lambda$ is its perspective composed with a linear map, which is convex on $\{t>0\}$. Since
$b/t$ is a convex combination of the fixed vectors $(h_{vy}-c)_y$, $\Lambda(z)=O(t)$, so $\Lambda$ extends continuously
by $0$ and is convex on the whole cube $[0,1]^X$. At a vertex $z=\mathbf 1_{X'}$ with $X'\ne\emptyset$ we have
$b/t=(g'_y-c)_y$, where $g'$ is the column density of $X'$, so
$\Lambda(\mathbf 1_{X'})=(\Phi_c(X',Y)/\Phi_c(X,Y))^{1/w}\le1$ by maximality. A convex function on the cube is at most
its largest vertex value, so $\Lambda\le1$, which says $\|b_+\|_s\le H\,t^{1-w/r}$. By H\"older's inequality,
\[
H^s\,\E_X[z_vV_v]=\E_Y\big[u_y^{s-1}b_y\big]\le\big\|u^{s-1}\big\|_{s/(s-1)}\,\|b_+\|_s=H^{s-1}\|b_+\|_s,
\]
using $u_y>0$ (Lemma~\ref{lem:columns}). This proves~\eqref{eq:support}. Taking $z=\mathbf 1-\epsilon e_v$ gives
$1-\epsilon V_v/N\le(1-\epsilon/N)^{1-w/r}$, and letting $\epsilon\downarrow0$ gives $V_v\ge1-w/r$.
\end{proof}

\begin{remark}
For $w\ge1$ one may use convexity of $\Lambda^w$ directly. For $0<w<1$ that function need not be convex, and taking the
$1/w$-th power is what makes the argument work. Of the $918$ controls in the certified bank, $648$ have $w<1$.
\end{remark}

\subsection{An integral that rules out failure of both steps}

We now define the number whose sign certifies a local step. The parameter $p$ is the target red density,
$\mu$ controls the blue step, and $x$ is the fraction of the potential that a red step must retain.
The parameters $w$ and $\beta$ are those of the potential.
Fix $0<p<1$, $0<\mu<1$, $w>0$, $0<\beta<1$ and $x>0$. For $q\in(0,\mu]$ let
\[
A(q)=\frac{1-q}{q}\log\frac{p(1-q)^w}{x},\qquad B(q)=w\log\frac\mu q,\qquad U(q)=e^{\beta B(q)}=(\mu/q)^{w\beta}.
\]
Let $M(q)=B(q)$ if $A(q)\le B(q)$, and otherwise
\[
M(q)=p\,A(q)+\frac{1-p}\beta\log\frac{U(q)-p\,e^{\beta A(q)}}{1-p}.
\]
The logarithm in the second branch is defined only when its numerator is positive.
Accordingly, the \emph{feasible set} $\feas=\feas(p,\mu,w,\beta,x)$ is the set of $q\in(0,\mu]$ with $A(q)\le B(q)$, or with
$A(q)>B(q)$ and $p\,e^{\beta A(q)}<U(q)$. When $\feas\ne\emptyset$ we define
\[
\cert(p,\mu,w,\beta,x)=\int_0^1\sup_{q\in\feas}\Big[q\,M(q)+w\,q\,(1+\log t)\Big]\,dt.
\]
On $\feas$ one has $M(q)\le B(q)$, by concavity of the logarithm, so the integrand is bounded above. The statement
$\cert<0$ includes the assertion that the integrand is integrable.

\paragraph{Meaning.} In the application, $q$ is the blue degree of a row of $X$, divided by $|X|$.
Passing to its blue neighbourhood retains a fraction $q^w$ of the factor $|X|^w$ in the potential.
The quantities $A(q)$ and $B(q)$ describe the changes in the column moment needed for the red and blue steps,
respectively. The function $M(q)$ bounds the average change allowed when both steps fail.
Row maximality gives a second, averaged constraint after we sort rows by blue degree.
Its contribution to the limiting inequality is $wq(1+\log t)$, where $t$ is the fraction of rows below
the corresponding degree. For each $t$, the supremum permits any feasible degree $q$.
If even this upper bound has negative integral, simultaneous failure is impossible.
Theorem~\ref{thm:envelope} makes this argument precise.

\subsection{The sharp envelope theorem}

Let $\mathcal B$ be any simple graph on $V$. For $i\in X$ let $B_i=\{j\in X:ij\in\mathcal B\}$, $q_i=|B_i|/N$, and let
$b_i(y)$ be the proportion of $B_i$ joined in red to $y$ (with $b_i(y)=0$ if $B_i=\emptyset$). In the application,
$\mathcal B$ is the blue graph and $B_i$ is the blue neighbourhood of $i$ in $X$.

We say row $i$ with $q_i>0$ \emph{fails blue} if
\begin{equation}\label{eq:B}
\sum_{y\in Y}(b_i(y)-c)_+^{\,s}<\Big(\frac\mu{q_i}\Big)^{w\beta}\sum_{y\in Y}(u_y)_+^{\,s},
\end{equation}
and row $i$ \emph{fails shifted-red} if
\begin{equation}\label{eq:R}
\sum_{y\in Y:\ h_{iy}=1}\Big((1-q_i)u_y-q_i\big(b_i(y)-g_y\big)\Big)_+^{\,s}
<(1-q_i)^{s-w\beta}\,x^\beta\,\pi_i^{\,1-\beta}\sum_{y\in Y}(u_y)_+^{\,s}.
\end{equation}

\begin{theorem}[sharp envelope]\label{thm:envelope}
Let $0<p<1$, $0<\mu<1$, $w>0$, $0<\beta<1$ and $0<x<\min\{p,(1-\mu)^w\}$. Suppose $\feas\ne\emptyset$ and
$\cert(p,\mu,w,\beta,x)<0$. Then there are $\eta,\tau,\zeta>0$ and $R_0$ with $R_0>w$ and $\beta R_0>1$ such that,
for every $r\ge R_0$, there is no configuration consisting of a host on $V$, a simple graph $\mathcal B$ on $V$,
disjoint nonempty $X,Y\subseteq V$ and a real $c$ with $|c-p|\le\eta$, such that
\begin{enumerate}[label=(\roman*)]
\item $g_y>c$ for every $y\in Y$;
\item $(X,Y)$ is maximal for $\Phi_c$ with parameters $w$, $r$, $s=\beta r$;
\item every $i\in X$ with $q_i>0$ fails blue;
\item there is $X_g\subseteq X$ with $|X\setminus X_g|\le\zeta|X|$ such that every $i\in X_g$ has $q_i\le\mu+\tau$
and fails shifted-red.
\end{enumerate}
\end{theorem}

In words: at a maximal pair whose shift is close to $p$, some row of positive blue degree passes the blue test, or more
than a $\zeta$-fraction of the rows either have large blue degree or pass the shifted-red test.

The proof is by contradiction. If both tests could fail for arbitrarily large $r$, we would obtain a sequence
of configurations with vanishing exceptional proportions. In the limit, maximality and the fact that the blue
graph is undirected give an averaged inequality incompatible with $\cert<0$.
We first state this limit argument, then recover the finite statement.

\begin{proposition}[limit form]\label{prop:limit}
Fix $p,\mu,w,\beta,x$ as in Theorem~\ref{thm:envelope}. Suppose that for each $j$ we have a configuration with
parameter $r_j$ satisfying (i)--(iv) of Theorem~\ref{thm:envelope} with $\tau=\tau_j$ and $\zeta=\zeta_j$, where
$r_j\to\infty$, $r_j>w$, $\beta r_j>1$, $\tau_j\to0$, $\zeta_j\to0$, the shifts $c_j$ and the moments $H_j$ are
bounded, and $\liminf_jd_j\ge p$ where $d_j=d_R(X_j,Y_j)$. Then, if $\feas\ne\emptyset$ and the integrand of $\cert$ is
integrable, $\cert(p,\mu,w,\beta,x)\ge0$.
\end{proposition}

\begin{proof}[Proof of Theorem~\ref{thm:envelope} from Proposition~\ref{prop:limit}]
Suppose the conclusion fails. Then for every $j\ge1$ the choice $\eta=\tau=\zeta=1/j$ and
$R_0=\max\{j,w+1,2/\beta\}$ admits some $r_j\ge R_0$ and a configuration with $|c_j-p|\le1/j$ satisfying (i)--(iv).
Then $|c_j|\le1+p$, $H_j\le\max_yu_y\le1-c_j\le2$, and $d_j=c_j+\E_Yu_y>c_j\ge p-1/j$. So
Proposition~\ref{prop:limit} applies and gives $\cert\ge0$, a contradiction.
\end{proof}

The proof has four main parts. We first bound what one row can contribute when both tests fail.
We then show that almost all rows have the same limiting red density.
Next we use row maximality and the symmetry of the blue graph to bound the average contribution from below.
Finally we pass to the limit and compare the result with the integral defining $\cert$.
We give all eight steps below.

The rest of this subsection proves Proposition~\ref{prop:limit}. Constants hidden in $O(\cdot)$ depend only on
$\beta$, $w$ and the bounds on $|c|$ and $H$. Fix one configuration. Since all $u_y>0$, put $a_y=u_y/H$, so $\E_Ya_y^s=1$,
and let $\nu$ be the probability measure on $Y$ with $\nu_y=a_y^s/M$. Put $m_i=\E_\nu h_{iy}$ and, for $q_i>0$,
$Z_i(y)=(b_i(y)-g_y)/u_y$. Since $u_y+(b_i(y)-g_y)=b_i(y)-c$,
\[
\E_\nu(1+Z_i)_+^s=\frac{\sum_y(b_i(y)-c)_+^s}{\sum_yu_y^s},
\]
\[
\E_\nu\Big[h_{iy}\big((1-q_i)-q_iZ_i\big)_+^s\Big]=\frac{\sum_{y:h_{iy}=1}\big((1-q_i)u_y-q_i(b_i(y)-g_y)\big)_+^s}{\sum_yu_y^s}.
\]
So~\eqref{eq:B} says $\E_\nu(1+Z_i)_+^s<U_i:=(\mu/q_i)^{w\beta}$, and~\eqref{eq:R} says
$\E_\nu[h_{iy}((1-q_i)-q_iZ_i)_+^s]<(1-q_i)^{s-w\beta}x^\beta\pi_i^{1-\beta}$. For $q_i=0$ the left side
of~\eqref{eq:R} is $m_i$.

\paragraph{Step 1: an optimisation on two parts of $Y$.} For a fixed row, split $Y$ according to whether its
edge to the row vertex is red. Jensen's inequality reduces the moment constraints to one number on each part.
Let $0<m<1$, $s>1$, $U>0$ and $t_U=U^{1/s}$. Consider the supremum of
$mz_{\rm in}+(1-m)z_{\rm out}$ subject to $z_{\rm in}\ge L$ and $m(1+z_{\rm in})_+^s+(1-m)(1+z_{\rm out})_+^s<U$. Raising
a coordinate below $-1$ to $-1$ keeps the constraint and raises the objective, so we may write
$\mathcal A=1+z_{\rm in}\ge0$ and $\mathcal C=1+z_{\rm out}\ge0$. By Jensen, $m\mathcal A+(1-m)\mathcal C<t_U$. If
$L\le t_U-1$, the supremum is $t_U-1$. The constraint is strict, so these boundary values are suprema, not maxima.
If $L>t_U-1$, then on the boundary of the closed constraint
$\mathcal C=((U-m\mathcal A^s)/(1-m))^{1/s}$ the objective has derivative $m[1-(\mathcal A/\mathcal C)^{s-1}]<0$ in
$\mathcal A$. Thus, when the strict constraint is feasible, the supremum is approached with $\mathcal A=1+L$
and $\mathcal C$ tending to its boundary value from below. It equals
\begin{equation}\label{eq:twocell}
mL+(1-m)\Big[\Big(\frac{U-m(1+L)^s}{1-m}\Big)^{1/s}-1\Big],
\end{equation}
The feasible set in this branch is nonempty exactly when $U>m(1+L)^s$.

For a row $i\in X_g$ with $q_i,m_i>0$, Jensen's inequality on the red cell $\{h_{iy}=1\}$ (of $\nu$-mass $m_i$),
applied to~\eqref{eq:R}, gives $m_i((1-q_i)-q_iz_{\rm in})_+^s<(1-q_i)^{s-w\beta}x^\beta\pi_i^{1-\beta}$, where
$z_{\rm in}=\E_\nu[Z_i\mid h_{iy}=1]$. Taking $s$-th roots, and using $(s-w\beta)/s=1-w/r$ and $\beta/s=1/r$,
\begin{equation}\label{eq:L}
z_{\rm in}>L_i:=\frac{(1-q_i)-(1-q_i)^{1-w/r}x^{1/r}(\pi_i^{1-\beta}/m_i)^{1/s}}{q_i},\qquad
q_irL_i=(1-q_i)\,r\big(1-e^{\ell_i/r}\big),
\end{equation}
with $\ell_i=\log x-w\log(1-q_i)+\beta^{-1}\big[(1-\beta)\log\pi_i-\log m_i\big]$. If $\pi_i,m_i\to d$ and
$q_i\to q$, then $rL_i\to A_d(q)$, where $A_d$ is $A$ with $p$ replaced by $d$. Jensen on both cells, applied
to~\eqref{eq:B}, gives the constraint of the two-cell problem with $U=U_i$, with $z_{\rm in}$ as above and
$z_{\rm out}=\E_\nu[Z_i\mid h_{iy}=0]$; its objective is $\E_\nu Z_i$.

\paragraph{Step 2: rows and columns concentrate.} Lemma~\ref{lem:columns} gives $a_y^s\ge1-\beta$. With
$K_\beta=-\log(1-\beta)/\beta$ and $\alpha=\E_Yu_y$, this gives $e^{-K_\beta/r}\le\alpha/H\le1$, where the upper bound is
the power-mean inequality. Maximality under restricting $Y$ to a subset $A$ of proportion $\theta$ gives
$\E_Y[a_y^s\mathbf 1_A]\le\theta^{1-s/r}$; applied to $A=\{a_y\ge t\}$ it gives $\Prob_Y(a_y\ge t)\le t^{-r}$ for
$t\ge1$. Let $J=\E_Ya^{s-1}\le1$ and $P_+=\E[(a-1)_+a^{s-1}]$. The layer-cake formula with the tail bound gives
\[
P_+\le\int_1^\infty\big(st^{s-1}-(s-1)t^{s-2}\big)t^{-r}\,dt=\frac r{(r-s)(r-s+1)}=O(1/r),
\]
and $1-J=\E[a^{s-1}(a-1)]\le P_+$.

Let $f_i=\E_Y[a^{s-1}h_{iy}]=HV_i+cJ$. Then $\E_Xf=H+cJ$. By Lemma~\ref{lem:rows}, $f_i-\E f\ge-Hw/r$, and since
$f-\E f$ has mean zero, $\E_X|f_i-\E f|\le2Hw/r$. Also $m_i=\E_Y[a^sh_{iy}]$, so
$|m_i-f_i|\le\E_Y[a^{s-1}|a-1|]=2P_+-(1-J)\le2P_+$. Finally $\E f-d=(H-\alpha)-c(1-J)$. Together,
\begin{equation}\label{eq:C1}
\E_X|m_i-d|=O(1/r).
\end{equation}
For the ordinary row densities, maximality and the power-mean inequality give
$\Prob_X(\pi_i\ge c+Ht)\le t^{-r/w}$ for $t>1$: the set $A$ of such rows has $H_c(A,Y)\ge Ht$, so
$(|A|/N)^wt^r\le1$. Integrating, $\E(\pi-c-H)_+\le Hw/(r-w)$, and since $\E\pi=c+\alpha$,
\begin{equation}\label{eq:C2}
\E_X|\pi_i-d|\le\frac{2Hw}{r-w}+2(H-\alpha)=O(1/r).
\end{equation}

\paragraph{Step 3: the averaged constraint from row maximality.} We sort rows by blue degree and apply
\eqref{eq:support} to initial portions of this ordering.
Let $\gamma=1-w/r$ and let $Q$ be the increasing quantile function of the law of $q_i$
under uniform $i\in X$, with distribution function $F(z)=\Prob(q\le z)$. Thus $Q(t)$ is the blue degree at
proportion $t$ in the ordered list, with the usual quantile interpretation at atoms. Integrating a function
of $Q(t)$ over $0<t<1$ is the same as averaging that function of $q_i$ over the rows.
The ordering lets us use row maximality to control every initial portion of the list, not just one row.
By the layer-cake formula and~\eqref{eq:support}
applied to $z_i=\mathbf 1[q_i\le z]$,
\[
\E[qV]=1-\int_0^1\E\big[V\mathbf 1_{q\le z}\big]\,dz\ge1-\int_0^1F(z)^\gamma dz=\int_0^1Q(t)\,\gamma t^{\gamma-1}\,dt.
\]
Subtracting $\E q=\int_0^1Q$,
\begin{equation}\label{eq:D1}
r\,\E[q(V-1)]\ge\int_0^1Q(t)\,k_r(t)\,dt,\qquad k_r(t)=r\big[(1-w/r)t^{-w/r}-1\big]\longrightarrow-w(1+\log t),
\end{equation}
with convergence in $L^1(0,1)$ (dominated by $C_wt^{-1/2}(1+|\log t|)$ for $r\ge2w$).

\paragraph{Step 4: symmetry of the blue graph.} Averaging over blue neighbours can be done in either order,
because $\mathcal B$ is undirected: $j\in B_i$ if and only if $i\in B_j$. Hence, for
each $y$, $\E_i[q_ib_i(y)]=N^{-2}\sum_i\sum_{j\in B_i}h_{jy}=\E_j[q_jh_{jy}]$. Since
$\sum_y\nu_y(h_{iy}-g_y)/u_y=V_i-1$, this gives
\begin{equation}\label{eq:ST}
\E_X\big[q_i\,\E_\nu Z_i\big]=\E_X\big[q_i(V_i-1)\big].
\end{equation}
Both~\eqref{eq:D1} and~\eqref{eq:ST} hold over all rows, including the exceptional ones.

\paragraph{Step 5: small blue degrees do not occur.} Call a row \emph{regular} if $|\pi_i-d|$ and $|m_i-d|$ are at most
$\eta_j$, where $\eta_j\to0$ slowly enough that, by~\eqref{eq:C1}, \eqref{eq:C2} and Markov's inequality,
the proportion of non-regular rows tends to $0$. Let $\kappa=\log(p/x)>0$. For regular rows in $X_g$ with
$0<q\le\delta_0$ and $j$ large, $\ell\le-\kappa/2$ and $m\ge p/2$, so~\eqref{eq:L} gives $q\,r\,L\ge\kappa/8$. On the
other hand, the blue constraint on the red cell gives $m(1+z_{\rm in})^s<U$, so
\[
q\,r\,z_{\rm in}<q\,r\big(e^{b/r}-1\big)\le q\,b\,e^{b/r},\qquad b=w\log(\mu/q)-\beta^{-1}\log m,
\]
and the right side is at most $C\sqrt q\,(1+\log(1/q))$ for $r\ge2w$, $m\ge p/2$ and $q\le\mu$. Since $z_{\rm in}>L$,
rows with $0<q<\delta$ cannot occur, for a fixed small $\delta>0$. At $q=0$, \eqref{eq:R} reads
$m<x^\beta\pi^{1-\beta}$, which is impossible when $\pi$ and $m$ are close to $d\ge p-o(1)>x$. Hence every weak limit
of the laws of $q_i$ is supported on $[\delta,\mu]$; the upper end comes from $q_i\le\mu+\tau_j$ on $X_g$.
Such weakly convergent subsequences exist because these are probability measures on the compact interval $[0,1]$.
The exceptional rows have vanishing proportion, so they contribute no mass outside $[\delta,\mu]$ in the limit.
This support restriction keeps the logarithms in the limiting argument away from $q=0$.

\paragraph{Step 6: the limiting envelope.} Pass to a subsequence with $d_j\to d_*\in[p,1]$.
The estimates in Step~2 make both $\pi_i$ and $m_i$ converge in probability to this same density.
The remaining variable is the blue degree $q$. We now find the largest scaled row gain compatible with a
given limiting value of $q$.
Combining
$z_{\rm in}>L$ with the blue constraint gives $m(1+L)_+^s<U$. At a regular limit point $q\in[\delta,\mu]$ we get
$rL\to A_{d_*}(q)$ and $(1+L)^s\to e^{\beta A_{d_*}(q)}$, so
\begin{equation}\label{eq:F}
d_*\,e^{\beta A_{d_*}(q)}\le U(q).
\end{equation}
The inequality is non-strict because a strict inequality can become equality in a limit.
In logarithms, \eqref{eq:F} says $T_{d_*}(q)\le\log x$, where
$T_d(q)=\big(1+\frac q{\beta(1-q)}\big)\log d+w\log(1-q)+\frac{wq}{1-q}\log\frac q\mu$. Since
$T_d'(q)=[\beta^{-1}\log d+w\log(q/\mu)]/(1-q)^2\le0$ for $d\le1$ and $q\le\mu$, we get $T_{d_*}(\mu)\le\log x$, that is,
\[
d_*\le\big[x/(1-\mu)^w\big]^{1/(1+\mu/(\beta(1-\mu)))}<1,
\]
using $x<(1-\mu)^w$. For $q$ in a compact subset of $(0,1)$ and $\pi,m\to d\in(0,1)$, $r$ times the two-cell
supremum~\eqref{eq:twocell} converges to $M_d(q)$, where $M_d$ is $M$ with $p$ replaced by $d$. In the non-binding case
$r(t_U-1)\to w\log(\mu/q)=B(q)$. In the binding case $m\,rL\to dA_d$ and
$r[((U-m(1+L)^s)/(1-m))^{1/s}-1]\to\beta^{-1}\log\frac{U-de^{\beta A_d}}{1-d}$. The two branches agree when $A_d=B$.
At the boundary $U=de^{\beta A_d}$ the scaled outside-cell term tends to $-\infty$, and outside the feasible set
$\feas_d$ (defined with $d$ in place of $p$) there are eventually no feasible parameters. So the limit is upper
semicontinuous with values in $[-\infty,\infty)$, where we set $M_d=-\infty$ off $\feas_d$.

\paragraph{Step 7: passage to the limit.} Jensen applied to~\eqref{eq:B} gives $(1+\E_\nu Z)^s<U$, so for $q>0$,
\[
q\,r\,\E_\nu Z\le q\,r\big[(\mu/q)^{w/r}-1\big],
\]
which is at most $w\mu^{w/r}q^{1-w/r}\log(\mu/q)$ for $q\le\mu$, uniformly bounded for $r\ge2w$, and is $\le0$ for
$q>\mu$. Rows with $q=0$ contribute $0$. So $q\,r\,\E_\nu Z$ has a common upper bound $H_0$ over all rows.
On regular rows in $X_g$ it is also at most $q$ times $r$ times the two-cell supremum.
Define a dominating random variable $W_j$ by taking the smaller of these two bounds on those rows, and $H_0$
on the exceptional rows. The latter have vanishing proportion.

Take a subsequence realising the $\limsup$ and then a weakly convergent subsequence of the joint laws of
$(q,m,\pi,\iota)$, where $\iota$ indicates that the row is regular and lies in $X_g$.
These laws live on a compact metric space. Skorokhod's theorem represents them on one probability space with
almost-sure convergence, while preserving their distributions and hence their expectations.
The limits have $m=\pi=d_*$ and $\iota=1$. In particular, almost every coupled row is non-exceptional for all
sufficiently large $j$, so Step~6 bounds the pointwise upper limit of $W_j$ by $qM_{d_*}(q)$.
Applying Fatou's lemma to the nonnegative variables $H_0-W_j$ now gives an upper bound on the limiting
expectation. Thus
\[
\limsup_j\E_X\big[q\,r\,\E_\nu Z\big]\le\int_0^1Q(t)\,M_{d_*}(Q(t))\,dt,
\]
where $Q$ is the quantile of the limit law of $q$. This integral is simply the expectation of the limiting
envelope, written in terms of quantiles.

For the lower bound we retain the averaged information from row maximality and blue-graph symmetry.
Equations~\eqref{eq:ST} and~\eqref{eq:D1} give
$\E_X[q\,r\,\E_\nu Z]=r\,\E[q(V-1)]\ge\int_0^1Q_jk_{r_j}$. The quantiles $Q_j$ converge to $Q$ at every continuity point
of $Q$ and are bounded by $1$, and $k_{r_j}\to-w(1+\log t)$ in $L^1$.
The monotone function $Q$ has at most countably many discontinuities. Dominated convergence handles the
product of $Q_j-Q$ with the limiting integrable kernel, while the $L^1$ convergence of the kernels handles
the remaining error. Therefore
\[
\liminf_j\E_X\big[q\,r\,\E_\nu Z\big]\ge-w\int_0^1Q(t)(1+\log t)\,dt.
\]
Together,
\begin{equation}\label{eq:N}
\int_0^1\Big[Q(t)\,M_{d_*}(Q(t))+w\,Q(t)(1+\log t)\Big]dt\ge0,
\end{equation}
and in particular the integral is finite. This is the constraint that will contradict a negative certificate:
the limiting degree distribution must satisfy it, even though no contradiction need hold for an individual row.

\paragraph{Step 8: from the limit to $\cert$.} Fix $q\in(0,\mu]$. On the binding branch put $E=e^{\beta A}$ and
$z=(U-dE)/(1-d)$. There $E>U>z>0$. Differentiating $M=dA+\frac{1-d}\beta\log z$ gives
$\beta\,\partial_dM=\log(E/z)+1-E/z\le0$ at fixed $A$ and $\partial_AM=d(1-E/z)<0$, while $\partial_dA_d=(1-q)/(qd)>0$.
So $M_d(q)$ decreases in $d$. On the non-binding branch $M_d=B$, and the branches meet continuously. Also
$de^{\beta A_d(q)}$ increases in $d$, so $\feas_d$ shrinks as $d$ grows. Since $d_*\ge p$, we get $\feas_{d_*}\subseteq\feas$
and $M_{d_*}\le M$ on $\feas_{d_*}$. Since~\eqref{eq:N} is finite, $Q(t)\in\feas_{d_*}$ for almost every $t$. For such
$t$,
\[
Q(t)M_{d_*}(Q(t))+wQ(t)(1+\log t)\le\sup_{q\in\feas}\big[qM(q)+wq(1+\log t)\big].
\]
Integrating gives $\cert\ge0$. This proves Proposition~\ref{prop:limit}, and hence Theorem~\ref{thm:envelope}. \qed

\begin{remark}
Two points of the proof deserve emphasis. A single row can fail both tests at the parameters we use; the contradiction
comes only from the averaged inequality~\eqref{eq:N}, that is, from the drift~\eqref{eq:D1} together with
stationarity~\eqref{eq:ST}. And Step~5 is essential: allowing $q=0$ in the supremum would make the integrand nonnegative
and destroy every certificate.
\end{remark}

%======================================================================
\section{Closing the induction inside one host}\label{sec:closure}
%======================================================================

This section turns the local step into a clique outcome for a pair $(X,Y)$ inside a host.
We assume an off-diagonal bound on proper subsets, so that we can stop when a subset is already large enough
to contain one of the required cliques. Otherwise we pass to a red neighbourhood, a blue neighbourhood,
or the pages of a blue book. Each move reduces a clique target and retains enough potential to continue.
The induction follows Lu and Wang~\cite[Lemma A.1]{LW}, with their density potential replaced by $\Phi_c$.

\begin{definition}[stopping rule]
For reals $C\ge1$, $\hat x,\hat y\in(0,1)$ and an integer $\ell\ge1$, a host on $V$ satisfies
$\mathrm{Stop}(C,\hat x,\hat y,\ell)$ if, for every integer $a\ge1$, every proper subset $Z\subsetneq V$ with
$|Z|\ge C\hat x^{-a}\hat y^{-\ell}$ contains a red $K_a$ or a blue $K_\ell$.
\end{definition}

\begin{theorem}[finite-host closure]\label{thm:closure}
Let $0<p<1$, $0<\mu<1$, $w>0$, $0<\beta<1$ and $0<x<\min\{p,(1-\mu)^w\}$, and suppose $\feas\ne\emptyset$ and
$\cert(p,\mu,w,\beta,x)<0$. Let $\xi,y_0,\nu,\hat x,\hat y\in(0,1)$ satisfy
\[
\xi<\min\{x,\hat x\},\qquad y_0<\hat y,\qquad \nu<\mu.
\]
Then there is an integer $L_0$, depending only on these ten numbers, with the following property. Let $C\ge1$, let
$\ell\ge L_0$, and let a host on $V$ satisfy $\mathrm{Stop}(C,\hat x,\hat y,\ell)$. Let $k,t\ge1$ be integers and
$X,Y\subseteq V$ disjoint and nonempty with
\[
d_R(X,Y)\ge p,\qquad |X|^w|Y|\ge C^{w+1}\xi^{-k}y_0^{-\ell}\nu^{-wt}.
\]
Then $X\cup Y$ contains a red $K_k$, or $X$ contains a blue $K_t$, or $Y$ contains a blue $K_\ell$.
\end{theorem}

The order of quantifiers matters later: $L_0$ is chosen before $C$ and before the host.

\subsection{Two lemmas on books}

\begin{lemma}[good pages]\label{lem:pages}
Let $(X,Y)$ be maximal with $H=H_c(X,Y)>0$, and let
\[
D_r=\frac w{r-w}+1-(1-\beta)^{1/(\beta r)},\qquad \mathcal G=\{v\in X:\pi_v\ge c\}.
\]
Then $1-|\mathcal G|/|X|\le D_r$.
\end{lemma}

\begin{proof}
For $t>1$, the rows with $\pi_v\ge c+Ht$ form a set $A$ whose column densities average at least $c+Ht$, so
$H_c(A,Y)\ge Ht$ by the power-mean inequality. Maximality gives $(|A|/N)^wt^r\le1$, so
$\Prob_X(\pi_v\ge c+Ht)\le t^{-r/w}$. The layer-cake formula gives $\E_X(\pi_v-c-H)_+\le Hw/(r-w)$. Since
$\E_X\pi_v=c+\alpha$ and $\alpha\ge(1-\beta)^{1/s}H$ (Lemma~\ref{lem:columns}),
\[
\E_X(c+H-\pi_v)_+=\E_X(\pi_v-c-H)_+-(\alpha-H)\le D_rH.
\]
A row outside $\mathcal G$ has $(c+H-\pi_v)_+>H$, so Markov's inequality gives the claim.
\end{proof}

\begin{lemma}[blue book with good pages]\label{lem:book}
Let $X$ be a nonempty vertex set of a host, $N=|X|$, and $\mathcal G\subseteq X$. Let $0<\theta<\tau<1$, $D$ real, and
$b\le m$ positive integers with
\[
1-\frac{|\mathcal G|}N\le D\le\frac14,\qquad \frac mN\le\frac14,\qquad D+\frac mN+\frac bm\le\tau-\theta.
\]
Let $W$ be the set of $v\in X$ with at least $\tau N$ blue neighbours in $X$. If $|W|\ge R(k,m)$, then $X$ contains a
red $K_k$, or there are disjoint $S,T\subseteq X$ such that $S$ is a blue clique of exactly $b$ vertices, all $S$--$T$
edges are blue, $T\subseteq\mathcal G$, and $|T|\ge\theta^bN/2$.
\end{lemma}

\begin{proof}
If $W$ has no red $K_k$, it has a blue clique $U$ with $|U|=m$. Let $P=\mathcal G\setminus U$; then
$|P|\ge(1-D)N-m\ge N/2$. Each $u\in U$ has at least $(\tau-D)N-m$ blue neighbours in $P$, so the blue density $\sigma$
between $U$ and $P$ satisfies $\sigma\ge\tau-D-m/N\ge\theta+b/m$. Choose a $b$-subset $S\subseteq U$ uniformly at random,
and let $T$ be the set of $v\in P$ joined in blue to all of $S$. If $v$ has $j_v$ blue neighbours in $U$, then
$\Prob(v\in T)=\binom{j_v}b/\binom mb$. The piecewise-linear interpolation of $j\mapsto\binom jb$ is nondecreasing and
convex, so by Jensen's inequality
\[
\E|T|\ge|P|\frac{\binom{\lfloor\sigma m\rfloor}b}{\binom mb}\ge|P|\Big(\sigma-\frac bm\Big)^b\ge\frac N2\,\theta^b,
\]
since each factor $(\lfloor\sigma m\rfloor-i)/(m-i)$ with $0\le i<b$ is at least $\sigma-b/m>0$. Some choice of $S$
achieves the expectation.
\end{proof}

\subsection{The local alternative}

For $v\in X$ write $B_v$ for the blue neighbourhood of $v$ in $X$, $q_v=|B_v|/|X|$, $R_v$ for its red neighbourhood in
$X$, and $Y_v$ for its red neighbourhood in $Y$.

\begin{lemma}\label{lem:local}
Fix an integer $N_0$ with $N_0(1-\mu)>1$. Under the hypotheses of Theorem~\ref{thm:closure} on $p,\mu,w,\beta,x$, there
are $\eta,\tau,\zeta>0$ and $R_0$ with $\eta<\min\{p,1-p\}$, $\mathrm{cap}:=\mu+\tau<1-1/N_0$, $R_0>w$ and $\beta R_0>1$,
such that the following holds for every $r\ge R_0$. Let $(X,Y)$ be maximal for $\Phi_c$ with $|c-p|\le\eta$,
$N=|X|\ge N_0$, $H>0$, and at most $\zeta N$ rows with $q_v>\mathrm{cap}$. Then for every
$\Delta\ge(1-c)/(N(1-\mathrm{cap})-1)$ and $c'=c-\Delta$,
\begin{itemize}
\item[(RED)] some $v$ with $q_v\le\mathrm{cap}$ has $\Phi_{c'}(R_v,Y_v)\ge x\,\Phi_c(X,Y)$, or
\item[(BLUE)] some $v$ has $\Phi_c(B_v,Y)\ge\mu^w\Phi_c(X,Y)$.
\end{itemize}
\end{lemma}

\begin{proof}
Take $\eta,\tau,\zeta,R_0$ from Theorem~\ref{thm:envelope} and shrink $\eta,\tau$ as required; shrinking keeps the
theorem true. Suppose both alternatives fail. Apply Theorem~\ref{thm:envelope} with $\mathcal B$ the blue graph. Its
condition (i) is Lemma~\ref{lem:columns}. Failure of (BLUE) at $v$ with $q_v>0$ reads $q_v^wH_c(B_v,Y)^r<\mu^wH^r$; raising
to the power $\beta$ gives~\eqref{eq:B}.

Now let $X_g=\{v:q_v\le\mathrm{cap}\}$ and fix $v\in X_g$. Put $\sigma_v=|R_v|/N=1-1/N-q_v>0$, and let $\rho_y$ be the
red density from $R_v$ to $y$. Counting red neighbours of $y\in Y_v$ in $X=\{v\}\cup B_v\cup R_v$ gives
$Ng_y=1+q_vNb_v(y)+\sigma_vN\rho_y$, so
\[
\sigma_v(\rho_y-c')=(1-q_v)u_y-q_v(b_v(y)-g_y)+\Big[\sigma_v\Delta-\frac{1-c}N\Big],
\]
and the bracket is nonnegative by the choice of $\Delta$. With $a_y=u_y/H$ and $J=\E_Ya_y^{s-1}>0$ we have
$\E_Y[h_{vy}a_y^{s-1}]=HV_v+cJ$, which is positive because $V_v\ge1-w/r>0$ (Lemma~\ref{lem:rows}) and $c\ge p-\eta>0$.
So $Y_v\ne\emptyset$.
Failure of (RED) at $v$, raised to the power $\beta$, gives
\[
\sum_{y\in Y_v}(\rho_y-c')_+^s<x^\beta\sigma_v^{-w\beta}\pi_v^{1-\beta}\sum_{y\in Y}u_y^s.
\]
Using the identity and $\sigma_v\le1-q_v$, this implies~\eqref{eq:R} at $v$. So all hypotheses of
Theorem~\ref{thm:envelope} hold, a contradiction.
\end{proof}

\subsection{Proof of Theorem~\ref{thm:closure}}

\paragraph{Parameters.} Choose $\lambda,y$ with $\xi<\lambda<\min\{x,\hat x\}$ and $y_0<y<\hat y$. Take
$N_0,\eta,\tau,\zeta,R_0,\mathrm{cap}$ from Lemma~\ref{lem:local}, and $\theta$ with $\mu<\theta<\mathrm{cap}$. Fix
$r\ge R_0$ so large that $D_r<\min\{1/4,(\mathrm{cap}-\theta)/2\}$. Fix $0<\delta<\min\{p/2,\eta\}$, and for $n\ge2$
let $c_n=p-\delta/n$ and $\Phi_n=\Phi_{c_n}$. Put
\[
\mathcal A=\min\{\lambda/\xi,(\mu/\nu)^w\}>1,\quad h=\mathrm{cap}-\theta-D_r>0,\quad
a_0=\min\Big\{\frac1w\log\frac{\hat x}\lambda,\frac1w\log\frac{\hat y}y,\log\frac1\mu\Big\}>0,
\]
and, for $n\ge4$,
\[
b_n=\max\Big\{1,\Big\lceil\frac{r\log(n^2/\delta)+w\log2}{w\log(\theta/\mu)}\Big\rceil\Big\},\qquad
m_n=\Big\lceil\frac{2b_n}h\Big\rceil.
\]
Both are $O(\log n)$, and $b_n\le m_n$. Let $\kappa_0=\inf_{n\ge2}\mathcal A^n/n^r>0$. For $2\le k\le n$ we have
$R(k,m)\le n^{m-1}$ by Erd\H{o}s--Szekeres, and $n^{m_n-1}=e^{O((\log n)^2)}$, so
\[
K_1=\sup_{n\ge4}e^{-a_0n}\max\Big\{N_0,\ 4m_n,\ \frac{2m_n}h,\ \frac{n^{m_n-1}}\zeta,\
\frac{1+n(n-1)/\delta}{1-\mathrm{cap}}\Big\}<\infty.
\]
Let $L_0$ be an integer with $(y/y_0)^{L_0}\delta^r\kappa_0\ge1$ and $e^{a_0L_0}\ge K_1$.

\paragraph{The induction.} Fix $C\ge1$, $\ell\ge L_0$ and a host satisfying $\mathrm{Stop}(C,\hat x,\hat y,\ell)$. We
prove by strong induction on $n=k+t$:
\begin{quote}
$(I_n)$ every disjoint nonempty pair $X,Y$ with $\Phi_n(X,Y)\ge C^{w+1}\lambda^{-k}y^{-\ell}\mu^{-wt}$ has an
\emph{outcome}: a red $K_k$ in $X\cup Y$, a blue $K_t$ in $X$, or a blue $K_\ell$ in $Y$.
\end{quote}
If $k=1$ or $t=1$, any vertex of $X$ is an outcome.

\emph{The theorem follows from $(I_n)$.} Given $X,Y$ as in the theorem, $H_{c_n}(X,Y)\ge d_R(X,Y)-c_n\ge\delta/n$ by the
power-mean inequality. Hence
\[
\frac{\Phi_n(X,Y)}{C^{w+1}\lambda^{-k}y^{-\ell}\mu^{-wt}}\ge\Big(\frac\delta n\Big)^r\Big(\frac\lambda\xi\Big)^k
\Big(\frac y{y_0}\Big)^\ell\Big(\frac\mu\nu\Big)^{wt}\ge\Big(\frac y{y_0}\Big)^\ell\delta^r\frac{\mathcal A^n}{n^r}\ge1.
\]

\emph{Inductive step} ($k,t\ge2$, so $n\ge4$). Among all pairs of nonempty subsets of $X$ and $Y$, choose one that
maximises $\Phi_n$. It is maximal, its potential is at least $\Phi_n(X,Y)$, and an outcome for it is an outcome for
$(X,Y)$; rename it $(X,Y)$. Write $N=|X|$, $M=|Y|$ and $H=H_{c_n}(X,Y)\in(0,1]$, so $\Phi_n\le N^wM$.

If $M\ge C\hat x^{-k}\hat y^{-\ell}$, then $Y$ is a proper subset of $V$, and the stopping rule gives a red $K_k$ or a
blue $K_\ell$ in $Y$. Otherwise $N^wM\ge C^{w+1}\lambda^{-k}y^{-\ell}\mu^{-wt}$ and $M<C\hat x^{-k}\hat y^{-\ell}$ give
\[
N>C\Big(\frac{\hat x}\lambda\Big)^{k/w}\Big(\frac{\hat y}y\Big)^{\ell/w}\mu^{-t}\ge e^{a_0(n+\ell)}\ge e^{a_0n}K_1.
\]
So $N\ge N_0$, $N\ge4m_n$, $N\ge2m_n/h$, $R(k,m_n)\le\zeta N$ and $N(1-\mathrm{cap})-1\ge n(n-1)/\delta$.

\emph{Book case}: at least $R(k,m_n)$ rows have $q_v\ge\mathrm{cap}$. Lemma~\ref{lem:pages} gives
$1-|\mathcal G|/N\le D_r$ for $\mathcal G=\{v:\pi_v\ge c_n\}$. Apply Lemma~\ref{lem:book} with $\tau=\mathrm{cap}$,
$b=b_n$, $m=m_n$; its conditions hold since $D_r+m_n/N+b_n/m_n\le D_r+h=\mathrm{cap}-\theta$. If it gives a red $K_k$
we are done. Otherwise it gives a blue $b_n$-clique $S$ and a set $T\subseteq\mathcal G$ of at least $\theta^{b_n}N/2$
pages. If $b_n\ge t$, then $S$ contains a blue $K_t$. If $b=b_n<t$, the columns of $T$ have average red density at least
$c_n$, so $H_{c_{n-b}}(T,Y)\ge c_n-c_{n-b}\ge\delta/n^2$, and
\[
\frac{\Phi_{n-b}(T,Y)}{\Phi_n(X,Y)}\ge2^{-w}\theta^{wb}(\delta/n^2)^r\ge\mu^{wb}
\]
by the choice of $b_n$. So $(I_{n-b})$ applies to $(T,Y)$ with targets $(k,t-b)$. A blue $K_{t-b}$ in $T$ together
with $S$ is a blue $K_t$ in $X$; the other outcomes are outcomes.

\emph{Local case}: fewer than $R(k,m_n)\le\zeta N$ rows have $q_v\ge\mathrm{cap}$. Apply Lemma~\ref{lem:local} with
$c=c_n$ and $\Delta=c_n-c_{n-1}=\delta/(n(n-1))$, which is large enough. In case (RED),
$\Phi_{n-1}(R_v,Y_v)\ge x\Phi_n(X,Y)\ge C^{w+1}\lambda^{-(k-1)}y^{-\ell}\mu^{-wt}$, since $x>\lambda$. By $(I_{n-1})$ with
targets $(k-1,t)$, a red $K_{k-1}$ in $R_v\cup Y_v$ together with $v$ is a red $K_k$; the other outcomes are outcomes.
In case (BLUE), lowering the shift to $c_{n-1}$ can only increase $H$, so
$\Phi_{n-1}(B_v,Y)\ge\mu^w\Phi_n(X,Y)\ge C^{w+1}\lambda^{-k}y^{-\ell}\mu^{-w(t-1)}$. By $(I_{n-1})$ with targets
$(k,t-1)$, a blue $K_{t-1}$ in $B_v$ together with $v$ is a blue $K_t$ in $X$. \qed

%======================================================================
\section{Profiles and banks}\label{sec:banks}
%======================================================================

Theorem~\ref{thm:closure} needs a stopping rule, that is, an off-diagonal Ramsey bound on proper subsets of the host.
This section shows how a finite collection of bounds can supply each other's stopping rules.
We call this collection a \emph{bank}. It contains bounds without a density restriction, called source pairs,
and bounds for hosts of a specified density, called profiles. Controls specify the parameters with which the
local step converts a source bound into a profile.

\paragraph{A two-bound example.}
Suppose we have two density-dependent bounds, with a common constant $C\ge1$.
The first says that a host of red density at least $1/2$ contains a red $K_k$ or a blue $K_b$ when
$|V|\ge Ce^{kL_R(b/k)}$. The second has the same form with a function $L_B$, and can be applied after
interchanging the colours and the targets. We use both bounds only at ratios in their respective domains.

Put $r=b/k$. Every host has red density at least $1/2$ or blue density at least $1/2$.
In the first case the required threshold is $Ce^{kL_R(r)}$.
In the second, interchanging the colours changes the targets to $(b,k)$, so the threshold is
\[
Ce^{bL_B(k/b)}=Ce^{k\,rL_B(1/r)}.
\]
Thus $|V|\ge Ce^{k\max(L_R(r),\,rL_B(1/r))}$ guarantees a red $K_k$ or a blue $K_b$ without any
density assumption. For example, at $b=2k$ the two exponents are $kL_R(2)$ and $2kL_B(1/2)$.
If $L_R(r)=a_R+b_Rr$ and $L_B(s)=a_B+b_Bs$ are positive affine functions on the intervals under consideration,
then the exponent per $k$ is simply
\[
\max\{a_R+b_Rr,\ b_B+a_Br\}.
\]
A line $A+Br$ with $A,B>0$ lying above this maximum supplies a source bound on that interval.
In the full construction, separate comparisons with the Lu--Wang input handle the remaining ratios.
This example explains both the reciprocal ratio in a colour reversal and the finite line comparisons below.

\paragraph{Why feedback is not circular.}
A source bound is used by the local step only on a proper subset of the current host.
Assume all bank claims have been proved for hosts with fewer vertices than $V$.
We first obtain the profiles that use the local step on $V$. Next we combine these profiles on $V$,
as in the example, and finally obtain the source bounds on $V$.
Consequently a source bound may help prove a profile which later helps prove that same source bound:
the first use is on a smaller host. The full argument below makes this induction simultaneous for all
claims and chooses one constant $C$ for the finite collection.

\subsection{Controls and profiles}

\begin{definition}[source pairs]
A \emph{bank} has finitely many \emph{source pairs} $(A_h,B_h)$ with $A_h,B_h>0$. Each pair stands for the claim
\[
(\mathrm{SOURCE}_h)\qquad |V|\ge C\,e^{A_hk+B_hb}\ \Longrightarrow\ V\supseteq(k,b)\qquad\text{for all integers }k,b\ge1.
\]
We write $T_h(\lambda)=A_h+B_h\lambda$, so that $A_hk+B_hb=k\,T_h(b/k)$.
\end{definition}

\begin{definition}[controls]
A \emph{control} is a tuple $\gamma=(p,\mu,w,\beta,x;\sigma;\xi,y_0,\nu)$ with $0<p,\mu,\beta<1$, $w>0$,
$0<x<\min\{p,(1-\mu)^w\}$, $\feas\ne\emptyset$, $\cert(p,\mu,w,\beta,x)<0$, a source index $\sigma$, and
$0<\xi<\min\{x,e^{-A_\sigma}\}$, $0<y_0<e^{-B_\sigma}$, $0<\nu<\mu$. We write $p_\gamma=p$, $w_\gamma=w$, and
$A_\gamma=-\log\xi$, $B_\gamma=-\log y_0$, $b_\gamma=-\log\nu$.
\end{definition}

A control packages everything Theorem~\ref{thm:closure} needs, with $\hat x=e^{-A_\sigma}$ and $\hat y=e^{-B_\sigma}$.
The source index says which source claim supplies the stopping rule.

\begin{definition}[profiles]
A \emph{profile} $P=(p_P,I_P,L_P)$ has a density $p_P\in(0,1)$, a compact interval $I_P\subset(0,\infty)$, and a
positive continuous piecewise-affine function $L_P$ on $I_P$ with finitely many rational knots. With constant $C\ge1$,
it stands for the claim
\[
(\mathrm{ASSERT}_P)\qquad d_R(V)\ge p_P\ \text{and}\ |V|\ge C\,e^{kL_P(b/k)}\ \Longrightarrow\ V\supseteq(k,b),
\]
for all integers $k,b\ge1$ with $b/k\in I_P$ and every host $V$.
\end{definition}

A profile is an exponential Ramsey bound for hosts of red density at least $p_P$, along the ratios $b/k\in I_P$. Note
that it is stated at integer pairs $(k,b)$.

\subsection{How profiles are justified}

A bank has two kinds of profiles. A \emph{direct} profile is obtained by applying the local step, possibly
after a sequence of blue steps. A \emph{combined} profile is obtained by combining direct bounds on the same
host, as in the two-bound example. The certificate labels these kinds RAW and CLOSED, respectively.
Both kinds are proved by the bank validity theorem; the labels describe their derivations, not their proof status.
A \emph{realization} is the data and inequalities that justify a profile on an interval of ratios.
We now describe the derivations used in the certificate.

\begin{definition}[direct realizations]
Each direct profile $P$ has a finite cover of $I_P$ by closed intervals $J$. Each $J$ carries a margin $\varepsilon_J>0$
and one of the following.
\begin{enumerate}[leftmargin=*]
\item \emph{Leaf.} A control $\gamma$ with $p_\gamma\le p_P$, such that on $J$
\[
L_P(r)\ge\varepsilon_J+\frac{A_\gamma+r\,(B_\gamma+w_\gamma b_\gamma)}{1+w_\gamma}.
\]
\item \emph{Retained route.} A control $\gamma$ with $\pi:=p_\gamma<p_P$; a partition
$\theta=u_0<u_1<\dots<u_M=\omega$ with $\theta>0$ and $J\subseteq[\theta,\omega]$; for each cell $J_j=[u_{j-1},u_j]$ a
slope $d_j$ and a \emph{child} profile $c(j)$ (any direct or combined profile of the bank) with $J_j\subseteq I_{c(j)}$; and
the continuous function $\psi$ on $[\theta,\omega]$ with $\psi(\theta)>0$ that is affine with slope $d_j$ on $J_j$.
These must satisfy
\begin{equation}\label{eq:route}
d_j>-\log\big(1-p_{c(j)}\big),\qquad \psi(s)>L_{c(j)}(s)\quad(s\in J_j),
\end{equation}
and, writing $A,B,b,w$ for $A_\gamma,B_\gamma,b_\gamma,w_\gamma$, on $J$
\[
L_P(r)\ge\varepsilon_J+\max\Big(\psi(r),\ \frac{A+rB+w\big(\theta b+\psi(r)-\psi(\theta)\big)}{1+w}\Big).
\]
\end{enumerate}
\end{definition}

A leaf applies Theorem~\ref{thm:closure} directly to a balanced cut. A route first spends blue steps along the ratio
interval $[\theta,r]$, at the cost recorded by $\psi$, and applies a child profile whenever the current set is dense
enough in red. Only when the blue target has dropped to $\theta k$ does it apply the local step. The density $\pi$ of
the terminal control is strictly below $p_P$; the gap pays for keeping only the rows with red density at least $\pi$.

\begin{definition}[combined realizations]
Each combined profile $Q$ has a finite cover of $I_Q$ by closed intervals $J$, each carrying one of the following.
Only direct profiles are used in these combinations, and the inequalities are non-strict.
\begin{enumerate}[leftmargin=*]
\item \emph{Density copy.} A direct profile $R$ with $p_R\le p_Q$, $J\subseteq I_R$ and $L_Q\ge L_R$ on $J$.
\item \emph{Complementary pair.} Direct profiles $R,S$ with $p_R+p_S\le1$, $J\subseteq I_R$, $\{1/r:r\in J\}\subseteq I_S$ and
$L_Q(r)\ge\max\big(L_R(r),\,r\,L_S(1/r)\big)$ on $J$.
\end{enumerate}
\end{definition}

\begin{definition}[source covers]
For each source pair $h$ fix $0<\ell_h<H_h$ and $\varepsilon_h,\eta_h>0$ with
\begin{equation}\label{eq:tail}
\fh(\lambda)\le T_h(\lambda)-\varepsilon_h(1+\lambda)\qquad(\lambda\in(0,\ell_h]\cup[H_h,\infty)),
\end{equation}
and a finite cover of $[\ell_h,H_h]$ by closed intervals $J$. Each $J$ is either \emph{grounded}, with
$\fh\le T_h-\eta_h(1+\lambda)$ on $J$, or carries two profiles $P_R,P_B$ with $p_{P_R}+p_{P_B}\le1$, $J\subseteq I_{P_R}$,
$\{1/\lambda:\lambda\in J\}\subseteq I_{P_B}$ and
\[
\max\big(L_{P_R}(\lambda),\,\lambda L_{P_B}(1/\lambda)\big)\le T_h(\lambda)-\eta_h(1+\lambda)\qquad(\lambda\in J).
\]
\end{definition}

The source covers are where the feedback happens: a source pair may be justified partly by the Lu--Wang function and
partly by the bank's own profiles.

\subsection{Validity of a bank}

\begin{theorem}[bank validity]\label{thm:bank}
Consider a bank whose direct profiles, combined profiles and source pairs have realizations and covers as above. Then there
is $C_*\ge1$ such that, with $C=C_*$, every claim $(\mathrm{ASSERT}_P)$ and every claim $(\mathrm{SOURCE}_h)$ holds for
every host.
\end{theorem}

The dependencies between profiles and sources may form cycles. The theorem is still true because every dependency is
used on a strictly smaller host. The proof is an induction on the number of vertices of the host.

For $C\ge1$ and a host $V$, let $\mathcal P_V(C)$ be the statement that every bank claim, with constant $C$, holds on
every proper subset $Z\subsetneq V$, under both namings of the colours.

\begin{lemma}[leaves and routes]\label{lem:leafroute}
For each leaf or route piece $J$ of a direct profile $P$ there are $K_J\ge1$ and $D_J\ge1$, depending only on the data of
the piece and not on $C$, with the following property. Let $V$ satisfy $\mathcal P_V(C)$, and fix a naming of the
colours. If $k\ge K_J$, $b/k\in J$, $d_R(V)\ge p_P$ and $|V|\ge D_JC\,e^{k(L_P(b/k)-\varepsilon_J/2)}$, then
$V\supseteq(k,b)$.
\end{lemma}

\begin{proof}
\emph{The stopping rule.} Let $\gamma$ be the control of the piece, with source index $\sigma$, and put
$\hat x=e^{-A_\sigma}$, $\hat y=e^{-B_\sigma}$. For every $W\subseteq V$ and every $\ell\ge1$, the host $W$ satisfies
$\mathrm{Stop}(C,\hat x,\hat y,\ell)$: a proper subset $Z\subsetneq W$ is a proper subset of $V$, and
$|Z|\ge C\hat x^{-a}\hat y^{-\ell}=Ce^{A_\sigma a+B_\sigma\ell}$ gives $Z\supseteq(a,\ell)$ by
$(\mathrm{SOURCE}_\sigma)$ on $Z$. The hypotheses of Theorem~\ref{thm:closure} on $(\xi,y_0,\nu,\hat x,\hat y)$ hold
by the definition of a control, so it gives a threshold $L_0(\gamma)$.

\emph{Leaves} ($D_J=3$). Choose $K_J$ with $\lfloor\alpha k\rfloor\ge\max(1,L_0(\gamma))$ for $k\ge K_J$, where
$\alpha=\min J$. Let $r=b/k$ and $L=L_P(r)-\varepsilon_J/2$. Since $|V|\ge3Ce^{kL}\ge3$, Lemma~\ref{lem:cut} gives a cut
$V=X\sqcup Y$ with $d_R(X,Y)\ge d_R(V)\ge p_P\ge p_\gamma$ and $|X|,|Y|\ge Ce^{kL}$. Then
\[
\log(|X|^w|Y|)\ge(1+w)\log C+(1+w)kL\ge(1+w)\log C+kA_\gamma+bB_\gamma+wb\,b_\gamma,
\]
by the leaf inequality. This is the size hypothesis of Theorem~\ref{thm:closure} with $t=\ell=b\ge L_0(\gamma)$. So
$X\cup Y$ has a red $K_k$, or $X$ or $Y$ has a blue $K_b$.

\emph{Routes.} Write $E(r)=\max\big(\psi(r),\,(A+rB+w(\theta b+\psi(r)-\psi(\theta)))/(1+w)\big)$ and
$\varepsilon=\varepsilon_J/2$, so $L_P(r)-\varepsilon_J/2\ge E(r)+\varepsilon=:z$. Let
$G=\psi+E(r)-\psi(r)+\varepsilon/2$ on $[\theta,r]$, a vertical shift of $\psi$. Then
\[
z-G(r)=\varepsilon/2,\qquad z+wG(\theta)-(A+rB+w\theta b)\ge(1+w/2)\varepsilon.
\]
Let $V_d=|d_1|+\sum_{j<M}|d_{j+1}-d_j|$, let $\delta_0>0$ be the minimum of $\psi-L_{c(j)}$ over the closed cells, let
$\eta_0=\min_j(1-p_{c(j)}-e^{-d_j})>0$, and let $\psi_{\min}=\min\psi>0$; all are positive by~\eqref{eq:route}. Put
$D_J=3(1-\pi)/(p_P-\pi)$, and choose $K_J$ so that for $k\ge K_J$: $\lfloor\theta k\rfloor\ge\max(1,L_0(\gamma))$,
$k\varepsilon/2>V_d$, $k\delta_0>V_d$, and $\eta_0e^{k\psi_{\min}-V_d}\ge1$.

Suppose $|V|\ge D_JCe^{kz}$. Take the cut $X_0,Y$ of Lemma~\ref{lem:cut} with $d_R(X_0,Y)\ge p_P$, and let $X$ be the
vertices of $X_0$ with at least $\pi|Y|$ red neighbours in $Y$. Counting red edges gives
$|X|\ge\frac{p_P-\pi}{1-\pi}|X_0|\ge Ce^{kz}$; also $|Y|\ge Ce^{kz}$, and every nonempty $Z\subseteq X$ has
$d_R(Z,Y)\ge\pi$. Let $\ell=b$ and $n_0=\lfloor\theta k\rfloor$. For $s\in(\theta,r]$ let $d(s)$ be the slope of $\psi$ on
the cell containing $s$ (the left cell at a knot). Define $\Gamma(t)=kG(\theta)$ for $t\le n_0$ and
$\Gamma(t)=kG(\theta)+\sum_{n_0<\tau\le t}d(\tau/k)$ for $n_0<t\le\ell$. Comparing the sum with the integral, one cell
boundary at a time, gives $|\Gamma(t)-kG(t/k)|\le V_d$ for $n_0<t\le\ell$, and $\Gamma(t)\ge k\psi_{\min}-V_d$
throughout.

We claim: for $t=1,\dots,\ell$ and every $Z\subseteq X$ with $|Z|\ge Ce^{\Gamma(t)}$, either $Z\cup Y$ has a red
$K_k$, or $Z$ has a blue $K_t$, or $Y$ has a blue $K_\ell$.

For $t\le n_0$, the terminal gap gives
$\log(|Z|^w|Y|)\ge(1+w)\log C+k(wG(\theta)+z)>(1+w)\log C+kA+\ell B+wtb$, which is the size hypothesis of
Theorem~\ref{thm:closure} for $(Z,Y)$ with the control $\gamma$ and $\ell\ge n_0\ge L_0(\gamma)$. For $n_0<t\le\ell$, let
$s=t/k$ and let $c$ be the child of the cell of $s$. Then $\Gamma(t)\ge kL_c(s)+k\delta_0-V_d>kL_c(s)$. If
$d_R(Z)\ge p_c$, then $(\mathrm{ASSERT}_c)$ on $Z$, which is a proper subset of $V$, gives $Z\supseteq(k,t)$. Otherwise
some vertex $v\in Z$ has more than $(1-p_c)(|Z|-1)\ge e^{-d(s)}|Z|$ blue neighbours $Z'$ in $Z$, using $|Z|\ge1/\eta_0$.
Then $|Z'|\ge Ce^{\Gamma(t-1)}$, and the claim for $t-1$ applied to $Z'$, with $v$ added to a blue $K_{t-1}$, finishes
the step.

Finally $\Gamma(\ell)\le kG(r)+V_d=kz-k\varepsilon/2+V_d<kz$, so $|X|\ge Ce^{\Gamma(\ell)}$, and the claim with $Z=X$,
$t=\ell=b$ gives $V\supseteq(k,b)$.
\end{proof}

\begin{lemma}[combined profiles on the same host]\label{lem:closed}
Suppose every direct claim holds on the host $V$, under both namings of the colours. Then every combined claim holds on $V$.
\end{lemma}

\begin{proof}
Let $r=b/k$ lie in a piece $J$ of $Q$, and let $|V|\ge Ce^{kL_Q(r)}$. For a density copy, $d_R(V)\ge p_Q\ge p_R$ and
$L_R(r)\le L_Q(r)$, so $(\mathrm{ASSERT}_R)$ applies. For a complementary pair, if $d_R(V)\ge p_R$ then
$(\mathrm{ASSERT}_R)$ applies at $(k,b)$. Otherwise $d_B(V)>1-p_R\ge p_S$; renaming the colours, $(\mathrm{ASSERT}_S)$ at
the pair $(b,k)$ has threshold $Ce^{bL_S(k/b)}=Ce^{k\,rL_S(1/r)}\le Ce^{kL_Q(r)}$. In both orientations the outcome is a
red $K_k$ or a blue $K_b$ in the original names. The complementary case does not use $d_R(V)\ge p_Q$.
\end{proof}

\begin{proof}[Proof of Theorem~\ref{thm:bank}]
\emph{Constants.} Let $K$ be an integer with $K\ge K_J$ for every leaf and route piece, and
$K\varepsilon_J/2\ge\log D_J$ for every direct piece. Let $M_I$ be the larger of $1$ and the largest right endpoint of all the
intervals $I_P$, and let $c_0=\min_h\min(A_h,B_h)>0$. Put
\[
C_{\rm prof}=\max\{R(k,b):1\le k<K,\ 1\le b\le M_I(K-1)\},\qquad
C_{\rm small}=\max_{0\le j\le K-2}j!\,(2/c_0)^j,
\]
and $C_{\rm tail}=C_\varepsilon$ from Theorem~\ref{thm:LW} with
$\varepsilon=\min_h\min(\varepsilon_h,\eta_h)$. Let $C=\max(1,C_{\rm prof},C_{\rm small},C_{\rm tail})$.

\emph{Grounding.} (a) For a profile $P$ and $k<K$ with $b/k\in I_P$, we have $b\le M_I(K-1)$, so
$|V|\ge Ce^{kL_P(b/k)}\ge C\ge R(k,b)$ gives $V\supseteq(k,b)$ whatever the density. (b) For a source pair $h$,
if $b/k\notin(\ell_h,H_h)$, then~\eqref{eq:tail} gives $\fh(k,b)+\varepsilon(k+b)\le k\,T_h(b/k)$, so
Theorem~\ref{thm:LW} proves $(\mathrm{SOURCE}_h)$. The same computation works on grounded cells. (c) If
$\min(k,b)<K$, write $j=\min(k,b)-1$ and $n=\max(k,b)$; then
$\binom{n+j-1}j\le(2n)^j\le j!(2/c_0)^je^{c_0n}\le C\,e^{A_hk+B_hb}$, so Erd\H{o}s--Szekeres proves
$(\mathrm{SOURCE}_h)$.

\emph{Induction.} Suppose some claim fails, and let $V$ be a host of least order on which some claim fails, under some
naming of the colours. Then $\mathcal P_V(C)$ holds.

\emph{Stage 1: direct claims on $V$.} Let $P$ be direct, $b/k\in I_P$, $d_R(V)\ge p_P$ and $|V|\ge Ce^{kL_P(b/k)}$. If $k<K$,
use grounding. Otherwise take the piece $J\ni b/k$. Since $k\varepsilon_J/2\ge\log D_J$,
$|V|\ge D_JCe^{k(L_P(b/k)-\varepsilon_J/2)}$, and Lemma~\ref{lem:leafroute} applies. The same holds with the colours
renamed.

\emph{Stage 2: combined claims on $V$.} By Stage~1 and Lemma~\ref{lem:closed}.

\emph{Stage 3: source claims on $V$.} By grounding we may assume $k,b\ge K$ and that $\lambda=b/k$ lies in a cell $J$
carrying profiles $P_R,P_B$. If $d_R(V)\ge p_{P_R}$, then $(\mathrm{ASSERT}_{P_R})$ on $V$ applies, since
$kL_{P_R}(\lambda)\le k\,T_h(\lambda)=A_hk+B_hb$. Otherwise $d_B(V)>1-p_{P_R}\ge p_{P_B}$, and $(\mathrm{ASSERT}_{P_B})$
with renamed colours at $(b,k)$ applies, since $b\,L_{P_B}(1/\lambda)=k\lambda L_{P_B}(1/\lambda)\le A_hk+B_hb$.

So every claim holds on $V$, a contradiction. Take $C_*=C$.
\end{proof}

\begin{corollary}[diagonal bound]\label{cor:diag}
Suppose a combined profile $Q$ has $1\in I_Q$, and the piece of $Q$ containing $1$ is a complementary pair. Then
$R(k,k)\le\lceil C_*e^{kL_Q(1)}\rceil$ for every $k\ge1$.
\end{corollary}

\begin{proof}
By Theorem~\ref{thm:bank} the direct claims hold on every host, so the proof of Lemma~\ref{lem:closed} applies to every
host at $(k,k)$. For a complementary pair it never uses the density hypothesis of $Q$.
\end{proof}

%======================================================================
\section{The certified bank and the proof of Theorem~\ref{thm:main}}\label{sec:certificate}
%======================================================================

\subsection{The bank}

The certified bank $\mathcal B^*$ is a finite file of rational numbers (Appendix~\ref{app:data}).
It records the controls, profiles, and source comparisons needed for Theorem~\ref{thm:bank}.
We found the data by numerical search, using the bounds produced by one application of the method as inputs
to further applications. The proof does not depend on the sequence of searches or on their numerical accuracy.
It depends only on the final file: no floating-point number enters the proof, and every inequality below
is certified with exact rational or outward-rounded interval arithmetic.
Table~\ref{tab:bank} lists its contents.

\begin{table}[ht]
\centering
\begin{tabular}{lr}
\toprule
object & number\\
\midrule
controls & 918\\
\quad $\beta=999/1000$; $p_\gamma=j/80-10^{-6}$ with $p_\gamma\in[0.0125,0.95]$ & \\
\quad $\mu\in[1/20,4/5]$; $w\in[1/8,4]$, with $w<1$ for $648$ controls & \\
source pairs & 2713\\
\quad of which justified by a single tangent of $\fh$ & 733\\
\quad of which justified by covers using profiles & 1980\\
Direct profiles (RAW) & 1902\\
\quad leaf pieces / route pieces & 9899 / 7284\\
\quad applications of child profiles in routes & 304\,695\\
Combined profiles (CLOSED, including the root) & 1003\\
\quad density-copy pieces / complementary-pair pieces & 9908 / 1686\\
\bottomrule
\end{tabular}
\caption{The certified bank $\mathcal B^*$.}\label{tab:bank}
\end{table}

The \emph{root} is a combined profile $Q^*$ with $p_{Q^*}=1/2$ and $I_{Q^*}=[1,81/80]$. It has a single piece, a
complementary pair using the same direct profile $R^*$ in both orientations, with $p_{R^*}=1/2$ and $I_{R^*}=[79/80,81/80]$.
Thus $p_{R^*}+p_{R^*}=1$, and
\[
L_{Q^*}(1)=L_{R^*}(1)=1.252369589401<z^*.
\]

\subsection{The checks}

The following properties of $\mathcal B^*$ were checked by an exact program. It uses rational arithmetic and
interval arithmetic with $224$-bit outward rounding for $\log$ and $\exp$; it rejects floating-point input.

\begin{enumerate}[label=(C\arabic*),leftmargin=*]
\item \emph{Controls.} For each of the $918$ controls: $0<x<\min\{p,(1-\mu)^w\}$; the point $q=\mu$ lies in the
feasible set $\feas$ (so $\feas\ne\emptyset$); and $\cert(p,\mu,w,\beta,x)<0$. The last check splits $(0,1)$ into
$8192$ cells, bounds the integrand above on each cell by a certified quantity, and verifies that the resulting upper bound
for the integral is negative. (The Lean version uses up to $65\,536$ cells for $151$ of the controls.) Also $\xi<x$, $\nu<\mu$, and
$-\log\xi>A_\sigma$, $-\log y_0>B_\sigma$ for the terminal source.
\item \emph{Direct pieces.} The leaf and route inequalities, with margins $\varepsilon_J>0$, the density conditions
$p_\gamma\le p_P$ for leaves and $p_\gamma<p_P$ for routes, the route conditions~\eqref{eq:route}, and the inclusions
$J_j\subseteq I_{c(j)}$.
\item \emph{Combined pieces.} The density copy and complementary pair inequalities, densities and domain inclusions,
including the reflected domains $\{1/r:r\in J\}$.
\item \emph{Sources.} The tail inequalities~\eqref{eq:tail} and the cover inequalities, with positive margins. Comparisons
with $\fh$ use tangent lines of $\fh$, which lie above $\fh$ by Lemma~\ref{lem:concave}, and enclosures of $f$ and $f'$
computed from the Lu--Wang data.
\item \emph{Root.} $Q^*$ is combined, $1\in I_{Q^*}$, its piece at $1$ is a complementary pair, $L_{Q^*}(1)<z^*$, and
$e^{z^*}<7/2$. The sharper comparison $e^{z^*}\le34986238/10000000$ is also kernel-checked
in the final Lean corollary.
\end{enumerate}

For the finite comparisons between profiles and tangent lines, a common refinement at their knots and at
the reciprocals of knots of reflected profiles makes each compared function affine on every cell.
Indeed, if $L(s)=a+bs$ on a cell, then $rL(1/r)=ar+b$ on the reflected cell.
Checking these comparisons at cell endpoints therefore proves them on whole intervals.
This argument concerns the profile and tangent-line comparisons, not the integral certificate or the original
Lu--Wang function $f$. The smallest margin in the source checks is $10^{-12}$;
combined inequalities are allowed to be tight, and the root inequality
$L_{Q^*}(1)\ge L_{R^*}(1)$ holds with equality.

\subsection{Proof of Theorem~\ref{thm:main}}

By (C1)--(C4), $\mathcal B^*$ satisfies all hypotheses of Theorem~\ref{thm:bank}.
In particular, every control satisfies the assumptions of Theorem~\ref{thm:closure}.
Theorem~\ref{thm:bank} gives a constant $C_*\ge1$.
By (C5) and Corollary~\ref{cor:diag},
\[
R(k,k)\le\big\lceil C_*e^{kL_{Q^*}(1)}\big\rceil
\le(C_*+1)e^{kL_{Q^*}(1)}\qquad(k\ge1).
\]
The second inequality uses $L_{Q^*}(1)>0$.
Set $\delta=z^*-L_{Q^*}(1)=0.000000100001>0$ and choose
\[
K=\max\left\{1,\left\lceil\frac{\log(C_*+1)}{\delta}\right\rceil\right\}.
\]
For $k\ge K$, we have $C_*+1\le e^{\delta k}$, and hence $R(k,k)\le e^{z^*k}=b_*^k$.
The certified comparison $e^{z^*}\le34986238/10000000<7/2$ gives the remaining eventual bounds.
For all $k\ge1$, taking $C=C_*+1$ gives $R(k,k)\le C e^{z^*k}$. \qed

%======================================================================
\section{Verification}\label{sec:verification}
%======================================================================

\subsection{Exact numerical checks}

The search for the certificate and the proof of its validity are separate.
The search uses floating-point arithmetic to find promising parameters.
The exact checker then reads the rational data and verifies (C1)--(C5), using outward-rounded enclosures
for logarithms and exponentials. Its integral bounds include every cell, including those near the endpoints.
It also checks the feasible-set witness $q=\mu$, all domain restrictions, and all positive margins required
by Theorem~\ref{thm:bank}. A serial run of the checker, without substituting recorded outputs, passed.
The Lean development independently proves the finite inequalities needed by the theorem.

\subsection{The Lean formalisation}

The whole argument has been formalised in Lean~4~\cite{Lean4} with Mathlib~\cite{Mathlib}.
This includes Theorems~\ref{thm:envelope},
\ref{thm:closure} and~\ref{thm:bank} and all lemmas used. It also includes the finite checks: the bank is translated
into Lean data, and every inequality of (C1)--(C5), including all $918$ integral certificates, is checked by the Lean
kernel through \texttt{decide}. We do not use \texttt{native\_decide} for these checks:
the kernel checks the proof terms rather than trusting the result of compiled evaluation.
The off-diagonal input, Theorem~\ref{thm:LW}, and the tangent lines of $\fh$ are taken from
Lu and Wang's Lean development~\cite{LWcode}. That development imports
\texttt{RamseyLean}~\cite{RamseyLean}, a Lean~4 formalisation of results from the CGMS argument.
These are mathematical dependencies, not merely sources of numerical data.

Related earlier formalisations include Mehta's Lean~3 proof~\cite{Mehta} and Paulson's Isabelle/HOL
proof~\cite{Paulson} of the exponential improvement of Campos, Griffiths, Morris and Sahasrabudhe.
We cite them as prior formal work; their source code is not imported into this Lean~4 proof.

\subsection{Scale and modular checking}

The main development uses Lean~4.33.1, while Lu and Wang's development uses Lean~4.32.1.
To obtain a checked chain in a single toolchain, we port the relevant modules to Lean~4.32.1.
In that chain every intermediate statement is supplied by its proof.
The pinned Mathlib revisions and reproduction instructions are given in Appendix~\ref{app:data}.

The publication archive contains $82\,190$ Lean source files with $4\,979\,699$ lines, occupying about
$1.82$\,GB before compression. This count excludes the external libraries.
Most of the source consists of generated certificate data and finite checks, rather than
handwritten mathematical proofs. The ported theorem closure contains $82\,143$ modules, supplemented by
$42$ bridge modules connecting the proof to the off-diagonal input.

We do not ask Lean to check all this source as one enormous expression or one compilation unit.
Each integral certificate is divided into cell computations. Separate modules prove the cell bounds,
and further modules combine them into bounds for the integral. The finite bank comparisons and the
off-diagonal input are also split into modules. Lean checks each module after the modules it imports.
The final modules combine these proved inequalities with the general mathematical theorems.

The count $82\,143$ refers to the ported local proof, not to every imported dependency.
The project inputs were also compiled from source in a separate clean build, with the pinned Mathlib and
its library dependencies supplied from their cache. We compared the resulting compiled modules with the
inputs used by the port. They agreed either byte-for-byte or, for five modules differing in stored source paths,
through matching hashes of all declaration types and proofs or values, with identical names and kinds.
The latter comparison is hash-based, not byte-for-byte, and the ported local proof was not rebuilt afterwards.
Thus verification used separate source builds and comparisons, not one monolithic rebuild.

Lean loads an imported compiled module without checking its proofs again.
The source compilations provide the proof checks; manifests and build stamps identify the sources and outputs
to which those checks apply. They do not replace proof checking.
Appendix~\ref{app:build} gives the complete build counts, comparisons and statement checks.

\subsection{Final statements and axiom checks}

In Lean, \lean{ramseyNumber k l} is the least $N$ such that every simple graph on $N$ vertices
contains a clique of size $k$ or an independent set of size $l$. A graph represents the red edges of a
colouring, so an independent set represents a blue clique.
The sharp theorem \lean{DiagRamsey.diagonal_le_exp_z} states that, for all sufficiently large $k$,
\[
(\texttt{ramseyNumber}\ k\ k:\R)
\le \exp\bigl((626184844701/500000000000:\R)\,k\bigr).
\]
The corollary \lean{DiagRamsey.diagonal_le_3p4986238_pow} states
\[
\exists\,K:\N,\ \forall\,k:\N,\ K\le k\ \Longrightarrow\
(\texttt{ramseyNumber}\ k\ k:\R)\le(34986238/10000000:\R)^k.
\]
Its proof includes the kernel-checked comparison $e^{z^*}\le34986238/10000000$.

The final axiom reports for the sharp theorem and the $3.4986238^k$ corollary list only
\[
\texttt{propext},\quad \texttt{Classical.choice},\quad \texttt{Quot.sound}.
\]
These are Lean's standard axioms for propositional extensionality, classical choice and quotient soundness.
The freshly compiled input theorems \lean{DiagRamsey.lu_wang_uniform_source_bound} and
\lean{DiagRamsey.lu_wang_lines_P0r12} also use only these three axioms.
No additional axiom, including \texttt{sorryAx}, occurs in these final theorem reports.

%======================================================================
\section{Remarks}
%======================================================================

\paragraph{Tight margins.} The certificate is exact, but its margins are small. The source checks have margins near
$10^{-12}$. The gap between the root profile exponent and $z^*$ is $0.000000100001$, and the difference
between $3.4986238$ and $e^{z^*}$ is about $6.2\times10^{-9}$.
A reimplementation must reproduce sufficiently precise enclosures of $f$; the source comparisons here use
precision better than $10^{-13}$.

\paragraph{Effectivity.} The constant $C_*$ and the threshold $K$ are not effective, because the constants of
Theorem~\ref{thm:envelope} come from a compactness argument.

\section*{Acknowledgements}

We thank Zhipeng Lu and Sichen Wang for making their mathematical argument, rational source data, exact checkers
and Lean development publicly available~\cite{LW,LWcode}. Their formalised off-diagonal bounds are an essential
input to our proof. We also thank the contributors to \texttt{RamseyLean}~\cite{RamseyLean}, which is imported
by that development.

We are grateful to the Lean developers~\cite{Lean4} and the Mathlib community~\cite{Mathlib} for the proof
assistant, mathematical library and proof automation on which our formalisation depends.
We also acknowledge the contributors to Batteries~\cite{Batteries} and Lean's standard library, which occur among the
dependencies. We thank Mehta and Paulson for their earlier formalisations of the exponential Ramsey
improvement~\cite{Mehta,Paulson}.

\appendix
\section{Data and reproduction}\label{app:data}

The paper, rational certificate, exact checker and Lean sources will be provided at \repository.
The following hashes identify the data and dependencies independently of that location.
Paths below are relative to the \texttt{diagonal-ramsey} directory.

\subsection{Data and pinned dependencies}

{\sloppy
\begin{itemize}[leftmargin=*,itemsep=2pt]
\small
\item Lu--Wang source data \lean{data/source.json}: SHA-256
\texttt{\seqsplit{f3e4a4eaf3f778ae125466eb6d47647a38fd38430fdcdffe8c5f774f9718c08c}}.
\item Lu--Wang repository commit: \texttt{\seqsplit{5ec98f470c43be747ea638f796d1494f9b4c0177}}.
\item Rational bank certificate: SHA-256 of the uncompressed file
\texttt{\seqsplit{bc3c24cf97cb9684b97f76ab7bf0e57775525d550a5dd1424d563900897442bc}}.
\item \texttt{RamseyLean} revision:
\texttt{\seqsplit{90e87da214701dd6eb3d56a2c7121839d8269d14}}.
\item Main development: Lean~4.33.1, Mathlib revision
\texttt{\seqsplit{0df444a360eaa60ab8c11dca51a86af692955474}}.
\item Ported development and Lu--Wang input: Lean~4.32.1, Mathlib revision
\texttt{\seqsplit{520045ab14e26149ee970e2e617ca04b09bde5d6}}.
\item Publication sources and build manifests: \lean{publication/P0r12/}.
The archive includes \lean{SharpBound.lean}, \lean{FinalCheck.lean} and the ported theorem closure.
\item SHA-256 of the recorded port build-output hash list:
\texttt{\seqsplit{dba4b8956cf2399570162edfe4759767a294592c50fae92b8b690da6263128e5}}.
\end{itemize}
\par}

\subsection{Reproducing the modular build}

The modular build instructions are in \lean{verify/README.md}; the entry point is
\lean{verify/verify.sh}. The default mode, \lean{MODE=single}, regenerates the certificate sources from the frozen
bank, checks the source manifest, and compiles the ported local proof under Lean~4.32.1, including all certificate,
bank and source-line modules. Here ``single'' means that each module is built in one toolchain,
not that the proof is checked as one large expression. The optional \lean{MODE=full} also builds the
Lean~4.33.1 development separately. Both modes can resume matching completed modules.
The setup can restore previously compiled input archives. To check these inputs from source, use a clean
\texttt{.lake} directory with the pinned dependencies, build the Lu--Wang, \texttt{RamseyLean}, bridge and
shared-definition modules, and compile the generated tangent-line modules separately.
The recorded clean build and its comparisons with the inputs used by the port are described below.
The publication archive supplies the sharp theorem sources and the instructions for checking them in the
ported environment. Both the script and the archive instructions should be read with their pinned toolchain
and dependency manifests; the two Lean versions are not interchangeable.

The build scripts schedule ready modules according to their memory requirements and preserve completed
build outputs. A stopped build can therefore resume without recomputing every certificate.
Generated sources are compared with a recorded manifest, and source/build stamps tie the saved outputs
to the sources from which they were checked.

\subsection{Recorded builds and comparisons}\label{app:build}

The publication validation matched all $82\,143$ closure stamps and the bridge-source manifests.
These hashes establish which sources and outputs were used; the mathematical verification comes from
the Lean kernel checks.

The separate clean input build compiled $770$ modules from source:
$741$ from the Lu--Wang development, $16$ from \texttt{RamseyLean}, $11$ bridge modules and two shared
definitions. It also compiled the $31$ tangent-line modules, giving $801$ input modules in total.
No archived project build outputs were used in this clean build.
The pinned Mathlib and its library dependencies were supplied from their cache.

We compared the resulting \texttt{.olean} files, which contain Lean's compiled modules,
with those used by the ported proof.
Of the $801$ files, $796$ were byte-identical. The other five differed in their stored absolute source-file paths.
For these five modules, a comparison of all $710$ declarations gave identical names and kinds,
with matching hashes of their types and of their proofs or values.
This last comparison is hash-based, not byte-for-byte.
The ported local proof was not rebuilt after this comparison.
Thus the verification uses separate source builds and comparisons of their outputs, rather than
one monolithic rebuild. Lean loads an imported \texttt{.olean} without checking its proofs again;
source hashes and build stamps alone would not replace the source compilation of these inputs.

The recorded port and publication checks also establish the following:
\begin{itemize}
\item all $82\,143$ ported modules compiled, with no unproved \lean{Theorems.*} statement modules remaining
in the final dependency closure;
\item the text of the $3.5^k$ statement matches its statement file, and \lean{Composite.lean} elaborates that
statement against the composed proof;
\item \lean{SharpBound.lean} compiles the sharp exponential theorem and its rounded corollaries, and
\lean{FinalCheck.lean} completes the publication check.
\end{itemize}
The final theorem and input axiom reports are stated in Section~\ref{sec:verification}.

\begin{thebibliography}{99}
\bibitem{CGMS} M.~Campos, S.~Griffiths, R.~Morris and J.~Sahasrabudhe, An exponential improvement for diagonal Ramsey,
\emph{Ann. of Math.} (2) 203 (2026), 869--932.
\url{https://doi.org/10.4007/annals.2026.203.3.4}.
\bibitem{ES} P.~Erd\H{o}s and G.~Szekeres, A combinatorial problem in geometry, \emph{Compositio Math.} 2 (1935),
463--470.
\bibitem{GNNW} P.~Gupta, N.~Ndiaye, S.~Norin and L.~Wei, Optimizing the CGMS upper bound on Ramsey numbers,
arXiv:2407.19026, 2024.
\bibitem{LW} Z.~Lu and S.~Wang, Retained-set descent for diagonal Ramsey numbers, arXiv:2609.14525, 2026.
\bibitem{LWcode} Z.~Lu and S.~Wang, \emph{Diagonal Ramsey numbers: source data, exact checkers and Lean proofs},
software repository, revision \texttt{5ec98f47},
\url{https://github.com/sichen-wang/diagonal-ramsey-numbers}.
\bibitem{RamseyLean} \texttt{RamseyLean} contributors, \emph{RamseyLean: a Lean~4 formalisation of the CGMS
Ramsey bounds}, software repository, revision \texttt{90e87da2},
\url{https://github.com/snorin239/RamseyLean}.
\bibitem{Mathlib} The mathlib Community, The Lean Mathematical Library, in \emph{Proceedings of the 9th ACM
SIGPLAN International Conference on Certified Programs and Proofs} (CPP 2020), 367--381.
\url{https://doi.org/10.1145/3372885.3373824}.
\bibitem{Lean4} L.~de Moura and S.~Ullrich, The Lean 4 Theorem Prover and Programming Language,
in \emph{Automated Deduction---CADE 28}, Lecture Notes in Computer Science 12699 (2021), 625--635.
\url{https://doi.org/10.1007/978-3-030-79876-5_37}.
\bibitem{Batteries} Batteries contributors, \emph{Batteries: an extended standard library for Lean},
software repository, \url{https://github.com/leanprover-community/batteries}.
\bibitem{Mehta} B.~Mehta, \emph{A formal proof of an exponentially better upper bound on Ramsey numbers},
Lean~3 development, 2023, \url{https://github.com/b-mehta/exponential-ramsey}.
\bibitem{Paulson} L.~C.~Paulson, An Exponential Improvement for Diagonal Ramsey,
\emph{Archive of Formal Proofs}, September 2024,
\url{https://isa-afp.org/entries/Diagonal_Ramsey.html}.
\end{thebibliography}

\end{document}
